\documentclass[11pt,reqno]{ijmart}

\usepackage[a4paper,margin=3.5cm]{geometry}
\usepackage[dvipsnames]{xcolor}
\usepackage[english]{babel}
\usepackage[babel]{microtype}

\usepackage{amsmath,amssymb,amsthm,mathtools,mathrsfs,esint}
\usepackage{mathalpha}
\usepackage{enumitem}

\usepackage[T1]{fontenc}
\usepackage[utf8]{inputenc}
\usepackage{stix} 

\usepackage[
  breaklinks=true,
  pagebackref=true,
  hypertexnames=false
]{hyperref}
\hypersetup{
  colorlinks=true,
  pdfpagemode=UseNone,
  citecolor=ForestGreen,
  linkcolor=MidnightBlue,
  urlcolor=RoyalBlue,
  pdfstartview=FitH,
  pdftitle={Existence of the critical two-dimensional empirical W1 constant},
  pdfauthor={Heng Ma},
  pdfsubject={Empirical transport, Gaussian fields, renormalization, and convex viscosity flow}
}
\usepackage[nameinlink,noabbrev]{cleveref}

\newtheorem{theorem}{Theorem}[section]
\newtheorem{lemma}[theorem]{Lemma}
\newtheorem{proposition}[theorem]{Proposition}
\newtheorem{corollary}[theorem]{Corollary}
\theoremstyle{definition}
\newtheorem{definition}[theorem]{Definition}
\newtheorem{remark}[theorem]{Remark}

\newcommand{\R}{\mathbb R}
\newcommand{\E}{\mathbb E}
\newcommand{\1}{\mathbf 1}
\newcommand{\cL}{\mathcal L}
\newcommand{\KR}{\mathrm{KR}}
\newcommand{\dist}{\operatorname{dist}}
\newcommand{\Lip}{\operatorname{Lip}}


% Keep short title / authors / page numbers,
% remove IJM branding only 
\fancyhead[LO]{}
\fancyhead[RE]{} 
% First page: remove journal name and DOI,
% keep centered page number at bottom
\fancypagestyle{firstpage}{%
  \fancyhf{}
  \cfoot{\thepage}
  \renewcommand{\headrulewidth}{0pt}
  \renewcommand{\footrulewidth}{0pt}
}


\title{Existence of the Critical Two-Dimensional Empirical \(W_1\) Constant}
\author{Heng Ma}
\address[Heng Ma]{Faculty of Data and Decision Sciences, Technion---Israel Institute of Technology, Haifa 32000, Israel}
\email{hengmamath(at)gmail(dot)com}
\urladdr{\url{https://hengmamath.github.io}}
\date{\today}

\begin{document}
\begin{abstract}
Let \(X_1,X_2,\ldots\) be independent uniform points in the unit square and let
\(\lambda\) denote Lebesgue probability measure.  We prove that the critical
two-dimensional empirical \(1\)-Wasserstein expectation has an asymptotic
constant:
\[
  \lim_{N\to\infty}
  \sqrt{\frac{N}{\log N}}\,
  \mathbb E W_1\!\left(\frac1N\sum_{i=1}^N\delta_{X_i},\lambda\right)
  =c_\star.
\]
The constant is characterized by a convex viscosity flow.  If \(U\) denotes the
unique solution in the convex cone-growth class of
\[
  \partial_tU=\frac1{4\pi}\sqrt{\det D^2U},
  \qquad U(0,q)=|q|,
\]
then \(c_\star=U(1,0)/\sqrt2\).  The proof replaces the empirical source
by a Gaussian field and successively absorbs groups of spatial scales
into a convex transport cost.  The quadratic response of this cost
determines the limiting flow. \textbf{The initial proof  was found by OpenAI's GPT-5.6 sol.}
\end{abstract}

\maketitle

\tableofcontents
\newpage 

\section{Introduction}\label{sec:introduction}

Let $X_1,X_2,\ldots$ be independent uniform points in $[0,1]^2$, let
$\lambda$ be Lebesgue probability measure on this square, and write
$\mu_N=N^{-1}\sum_{i=1}^N\delta_{X_i}$.  We study the leading constant
in the expected transport distance from $\mu_N$ to $\lambda$.
The Ajtai--Koml\'os--Tusn\'ady matching theory and Talagrand's
general-measure formulation give the critical estimate
\[
  \mathbb E W_1\!\left(\sum_{i=1}^N\delta_{X_i},N\lambda\right)
  \asymp \sqrt{N\log N};
\]
see Talagrand~\cite{Talagrand1992}.  This determines the order of the
expectation but leaves open whether its ratio to $\sqrt{N\log N}$
converges.  We prove convergence and characterize the limit by a
deterministic evolution equation.  The theorem is stated in the equivalent
empirical-process formulation; the equivalence with $N\mathbb E
W_1(\mu_N,\lambda)$ is recorded in Lemma~\ref{lem:empirical-transport}.

\begin{theorem}[Main result]\label{thm:intro-main}
Let \(X_1,X_2,\ldots\) be independent random variables uniformly distributed
on \([0,1]^2\), and let \(\lambda\) be Lebesgue probability measure on this
square.  There is a finite positive constant \(c_\star\) such that
\[
\lim_{N\to\infty}\frac1{\sqrt{N\log N}}
\mathbb E\sup_{\Lip(f)\leq1}
\left|\sum_{i=1}^N\left(f(X_i)-\int_{[0,1]^2}f\,d\lambda\right)\right|
=c_\star.
\]
Moreover, \(c_\star=\kappa/\sqrt2\), where \(\kappa=U(1,0)\) and \(U\) is
the unique convex cone-growth viscosity solution, in the sense specified
in Subsection~\ref{subsec:solution-class}, of
\[
\partial_tU=\frac1{4\pi}\sqrt{\det D_q^2U},
\qquad U(0,q)=|q|.
\]
\end{theorem}

The result concerns uniform sampling on the square with the Euclidean
transport cost.  It characterizes $c_\star$ without evaluating it
numerically.  We next explain why the transport problem leads to this
equation, before specifying its solution class.

\subsection{Proof strategy}\label{subsec:proof-roadmap}

Write $A_N$ for the expectation in Theorem~\ref{thm:intro-main}.
The proof uses the flux formulation of transport.  For a balanced source
$\sigma$ on a square $Q$, it reads
\[
\|\sigma\|_{\KR}
=\inf_{-\operatorname{div}v=\sigma,\ v\cdot n=0}|v|(Q).
\]
For an ordinary vector field, $|v|(Q)=\int_Q|v(x)|\,dx$.
The divergence constraint specifies the source transported by $v$, and
the boundary condition prevents flux from leaving $Q$.
Proposition~\ref{def:kr-beckmann} gives the precise formulation for vector
measures and its dual formulation with Lipschitz potentials.

Fluxes can be added when sources are added, but the cost of their sum
depends on their directions.  Thus separate optimal costs for different
scales do not determine the optimal cost of the full source.  We preserve
this dependence by keeping a convex function of the coarse flux as the
finer scales are removed.  Its evolution is the main object of the proof.

\smallskip
\paragraph{\textit{From empirical sampling to a Gaussian field.}}
For a $K\times K$ grid in the unit square, let $c_z$ be the cell centers,
let $G_z$ be independent standard Gaussians, and set
\[
\Gamma_K=\mathbb E\left\|
\sum_z(G_z-\bar G)\delta_{c_z}\right\|_{\KR},
\qquad \bar G=K^{-2}\sum_zG_z.
\]
Poissonization makes the cell counts independent.  Coupling these counts
to Gaussians and transporting the balanced coupling error gives
Lemma~\ref{lem:gaussianization-dyadic}:
\[
\left|A_N-\sqrt{N/K^2}\,\Gamma_K\right|
\leq C\left(\sqrt N+\frac NK+K\sqrt{\log(2K)}\right).
\]
For a dyadic choice $K_N\asymp\sqrt N/(\log N)^a$, with $0<a<1/2$,
these errors are negligible compared with $\sqrt{N\log N}$.
Thus it suffices to prove
\begin{equation}
\frac{\Gamma_{2^J}}{2^J\sqrt{J\log2}}\longrightarrow\kappa.
\tag{R}\label{eq:roadmap-gaussian-target}
\end{equation}
The relation $\log K_N/\log N\to1/2$ then accounts for the factor
$1/\sqrt2$ in the empirical constant.

\smallskip
\paragraph{\textit{Removing fine scales changes the cost of the coarse flux.}}
After dilation to $Q_{2^J}=[0,2^J]^2$ and spreading each Gaussian mass
over its unit cell, we obtain a balanced density $\rho_{2^J}$.
It is a sum of $J$ independent Haar details, each recording fluctuations
between dyadic children and their parent.  The error in spreading the
masses is $O(2^{2J})$, so the Gaussian target becomes
\[
\frac{\mathbb E\|\rho_{2^J}\|_{\KR}}{2^{2J}\sqrt{J\log2}}
\longrightarrow\kappa.
\]

Suppose that the current cost of a flux $p\in\mathbb R^2$ is $\Psi(p)$.
On a region where the coarse flux is almost constant, a fine source is
transported by a correction $u$, and the cost is applied to $p+u$.
Optimizing the fine correction and averaging over the fine source
therefore produces a new function of $p$, rather than a single number.
We denote this update by $\mathcal H_{h,L}\Psi$, where $h$ is the number
of fine levels being removed and $L$ is an auxiliary cell period.
Definition~\ref{def:cell-operator} constructs the update through the dual
potential problem with compatible boundary values.

The quantitative assertion behind this interpretation is
\[
\mathbb E J_\Psi^Q(\eta+\zeta)
\approx\mathbb E J_{\mathcal H_{h,L}\Psi}^Q(\zeta).
\]
Here $J_\Psi^Q$ is the transport cost with integrand $\Psi$, $\eta$ is the
fine source, and $\zeta$ is the independent coarse source.  The symbol
$\approx$ denotes the explicit error in
Lemma~\ref{lem:two-scale-scalarization}, not an exact identity.
Starting from the Euclidean cost, repeated replacement gives
\[
\Phi_0(p)=|p|,\qquad \Phi_{k+1}=\mathcal H_{h,L}\Phi_k.
\]
After the last source is removed, the remaining flux has zero divergence
and zero normal trace, hence zero mean.  Convexity then makes its minimum
cost per unit area exactly $\Phi_m(0)$.  This explains why the original
random optimization can be reduced to evaluating a deterministic
iteration at the origin.

\smallskip
\paragraph{\textit{Scale separation and common boundary values justify the replacement.}}
Two issues must be controlled in the preceding reduction.  The coarse
flux must vary little on a fine cell, and independently optimized cell
interiors must fit together.  To obtain the required separation, we keep
groups of $h$ consecutive levels and omit $b$ levels between groups.
Section~\ref{sec:dyadic} constructs a centered linear flux for each Haar
detail, with zero normal trace and uniformly bounded pointwise second
moment.  Independence gives an expected cost at most
$C2^{2J}\sqrt D$ for any $D$ omitted levels.  Consequently, omitting
$D=o(J)$ levels does not affect the desired limit.

For the boundary issue, prescribing a mean slope alone does not specify
the values of a potential along a shared face.  The cell operator retains
the full boundary function, including the compatibility of opposite
faces.  It chooses that boundary function after taking expectation, while
the optimized interior may depend on the fine source.  The original
optimizer can choose its boundary values using the fine source as well.
Lemma~\ref{lem:finite-trace-net} controls this additional freedom by a
finite approximation of the boundary functions and Gaussian concentration.
The two-scale proof then smooths and freezes the coarse optimizer, inserts
the cell optimizers, and joins them with an explicit error.  The primal
flux formulation gives the upper bound and the dual potential formulation
gives the lower bound.

\smallskip
\paragraph{\textit{The quadratic response determines the equation.}}
To identify the evolution of $\Phi_k$, consider first a smooth convex cost
near a background flux, with positive-definite Hessian $H$.  The
quadratic model for a unit outer block is
\[
\mathbb E\inf_{-\operatorname{div}u=\eta_h,\ u\cdot n=0}
\frac12\int_{Q_1}u^THu
=\frac12\operatorname{Tr}\!\left(\Pi_h(L_H^N)^{-1}\right),
\qquad L_H=-\operatorname{div}(H^{-1}\nabla).
\]
Here $\eta_h$ is the normalized $h$-level Gaussian band, $\Pi_h$ is its
covariance projection, and $L_H^N$ has Neumann boundary conditions.
The inverse matrix $H^{-1}$ appears because the dual quadratic cost is
$\frac12\nabla f^TH^{-1}\nabla f$.
In two dimensions the inverse elliptic eigenvalues have a logarithmic
sum.  The covariance projection cuts this sum off at spatial resolution
$2^{-h}$, and the elliptic eigenvalue count contributes
$\sqrt{\det H}$.  Lemma~\ref{lem:grouped-quadratic-generator} computes
the coefficient:
\[
\frac12\operatorname{Tr}\!\left(\Pi_h(L_H^N)^{-1}\right)
=\frac{h\log2}{4\pi}\sqrt{\det H}+o(h).
\]
The Dirichlet quadratic problem has the same leading term.  These two
problems bound the quadratic cell value, so the coefficient is independent
of the auxiliary period $L$.

The suggested generator is $\mathcal G(H)=(4\pi)^{-1}\sqrt{\det H}$,
with a time increment $h\log2$ per group.  To justify this evolution,
Section~\ref{sec:consistency} derives one-sided estimates
on smooth tests.  The upper estimate uses a Neumann flux.  The lower
estimate requires a Dirichlet potential whose gradient respects the
Lipschitz constraint, followed by a global convex barrier that makes a
local contact usable in the cell operator.  These estimates are formulated
for viscosity tests because the iterated costs need not be smooth.

\smallskip
\paragraph{\textit{Diffusive rescaling and the accumulated error.}}
The stability estimate $|p|\leq\Phi_k(p)\leq|p|+C\sqrt{kh}$ suggests
rescaling both the cost and its flux variable by the square root of time.
For a large time scale $r$, the relevant function is
\[
r^{-1/2}\Phi_{\lfloor sr/(h\log2)\rfloor}(\sqrt r\,q).
\]
For the spatial tests $\psi$ covered by Section~\ref{sec:consistency},
the rescaled one-sided estimates have leading term
$(h\log2/r)\mathcal G(D_q^2\psi)$, with time step $h\log2/r$.
In the joint regime $h\to\infty$, $L\to\infty$, and $h/r\to0$,
Lemma~\ref{lem:viscosity-limit} identifies its limit with $U(s,q)$.
The consistency estimates give the equation; comparison in the convex
class makes the limit unique.  In particular, at the retained time
$t_J=mh\log2\sim J\log2$, this yields
$\Phi_m(0)/\sqrt{t_J}\to U(1,0)$.

The scale choices must make this limit compatible with all earlier
replacements.  If $e_J$ is the error per unit area in one replacement,
the required conditions are
\[
D=o(J),\qquad m e_J=o(\sqrt J),\qquad h/t_J\longrightarrow0.
\]
Section~\ref{sec:conclusion} uses $h=\lfloor J^{3/4}\rfloor$ and
$b=\lceil(\log J)^2\rceil$, together with intermediate geometric scales,
to verify these conditions in one sequence.  Thus
\[
2^{-2J}\mathbb E\|\rho_{2^J}\|_{\KR}
=\Phi_m(0)+o(\sqrt J),
\]
which gives \eqref{eq:roadmap-gaussian-target} and then the empirical
limit.

\subsection{Viscosity formulation and notation}\label{subsec:solution-class}

A convex cone-growth solution is a continuous function
\(U:[0,\infty)\times\mathbb R^2\to\mathbb R\) whose spatial slices are convex
and \(1\)-Lipschitz and which is nondecreasing in time.  On every finite time
interval, \(U(s,q)-|q|\) is bounded uniformly in \(q\), and \(U\) has a
spatially uniform time modulus.  We use the strong convex viscosity
formulation: every smooth upper test satisfies the subsolution inequality,
and every smooth lower test with positive-semidefinite spatial Hessian
satisfies the supersolution inequality, including when that Hessian is
singular.  Lemma~\ref{lem:viscosity-limit} records the precise formulation and
proves comparison in this class.

\smallskip
\paragraph{\textit{Notation.}}
We write \(Q_L=[0,L]^2\), so \(|Q_L|=L^2\), and all vector norms are
Euclidean.  Constants denoted by \(C\) may change from line to line.  A vector
measure on a square has zero normal trace when its extension by zero has no
boundary divergence measure; this convention is made precise in
Proposition~\ref{def:kr-beckmann}.  A source is \emph{balanced} on a square if
its total mass there is zero.  We write \(B_r(a)\) for the Euclidean ball of
radius \(r\) centered at \(a\), set \(B_r:=B_r(0)\), and write
\(\mathbb S_+^2\) for the cone of symmetric
positive-semidefinite \(2\times2\) matrices.

\section{Transport formulation and Gaussian reduction}\label{sec:transport}

We first express the empirical process as a transport cost and then compare
it with a Gaussian grid.  The comparison keeps track of the grid size, so
the eventual Gaussian limit will imply convergence for every integer $N$.

\subsection{Transport costs and fluxes}

\begin{proposition}[Transport costs and duality]\label{def:kr-beckmann}
For a finite signed measure $\sigma$ of total mass zero on a bounded convex
set $Q\subset\mathbb R^2$, put
\begin{equation*}
\|\sigma\|_ {\rm KR}
:=\sup_ {\operatorname {Lip}(f)\leq1}\int_ Q f\,d\sigma .
\end{equation*}
Let $\mathcal M(Q;\mathbb R^2)$ be the vector Radon measures on
$Q$.  For $v\in\mathcal M(Q;\mathbb R^2)$, denote its extension by
zero to $\mathbb R^2$ by $\widetilde v$.  We write
\begin{equation*}
-\operatorname {div}v=\sigma,\qquad v\cdot n=0,
\end{equation*}
to mean the single distributional identity
$-\operatorname {div}\widetilde v=\widetilde\sigma$ on
$\mathbb R^2$.  Thus no boundary measure is present; this is the
definition of zero normal trace used below.  Then
\begin{equation}
\|\sigma\|_ {\rm KR}
=\inf\left\{|v|(Q):
-\operatorname {div}v=\sigma,\quad v\cdot n=0\text{ on }\partial Q
\right\}.
\tag{1}
\end{equation}
More generally, for a convex function $\Psi:\mathbb R^2\to\mathbb R$
with recession function $|\cdot|$, define
\begin{equation*}
J^Q_ \Psi(\sigma)
:=\inf_ {-\operatorname {div}v=\sigma,\ v\cdot n=0}
\left\{\int_ Q\Psi(v^a(x))\,dx+|v^s|(Q)\right\},
\end{equation*}
where $v=v^a\,dx+v^s$ is the Lebesgue decomposition.  This is the
relaxation of the integral with linear growth.
Its dual form is
\begin{equation}
J^Q_ \Psi(\sigma)
=\sup_ {f\in C^1(\overline Q)}
\left\{\int_ Q f\,d\sigma-\int_ Q\Psi^*(\nabla f)\right\}.
\tag{2}
\end{equation}
The same value is obtained with $W^{1,\infty}(Q)$ potentials.

\end{proposition}

\begin{proof}
We prove (2) by separating the relaxed primal epigraph and then take
$\Psi(z)=|z|$ to obtain (1).  Working with vector measures retains
singular fluxes and preserves the zero-normal-trace constraint under
weak-* limits.  The required closedness is verified next.

Set
\begin{equation*}
I(v)=\int_ Q\Psi(v^a)\,dx+|v^s|(Q).
\end{equation*}
The epigraph
\begin{equation*}
\mathscr E=\{(-\operatorname {div}\widetilde v,r):
v\in\mathcal M(Q;\mathbb R^2),\ I(v)\leq r\}
\end{equation*}
is convex and closed for distributional convergence in the first
coordinate and ordinary convergence in the second.  Indeed a bounded
energy sequence is weak-* compact as measures, the divergence constraint
passes to the limit, and the displayed relaxed integral is weak-* lower
semicontinuous (this last assertion follows directly by writing it as
the supremum of $\int a\cdot dv-\int\Psi^*(a)$ over continuous
$a$).

Separate $(\widetilde\sigma,r)$, with $r<\inf I$, from
$\mathscr E$.  The separating distribution is a smooth test function
$f$, and the coefficient of the epigraph coordinate is positive;
normalize it to one.  Minimizing the resulting affine functional over
absolutely continuous and singular parts of $v$ gives respectively
$-\int\Psi^*(\nabla f)$ and the constraint $|\nabla f|\leq1$.
This proves the reverse inequality in (2); integration by parts gives
the forward inequality.  Approximation of Lipschitz functions by smooth
ones on a convex square gives the $W^{1,\infty}$ version.  Taking
$\Psi(z)=|z|$, for which $\Psi^*$ is the indicator of the closed unit
ball, proves (1).

\end{proof}

\begin{lemma}[The empirical process as a transport cost]\label{lem:empirical-transport}
Let
\begin{equation*}
\nu_ N=\sum_ {i=1}^N\delta_ {X_ i},\qquad \sigma_ N=\nu_ N-N\lambda,
\end{equation*}
and let $A_ N$ denote the expectation in Theorem~\ref{thm:intro-main}. Then
\begin{equation}
A_ N=\mathbb E\|\sigma_ N\|_ {\rm KR}
=\mathbb E W_ 1(\nu_ N,N\lambda).
\tag{3}
\end{equation}

\end{lemma}

\begin{proof}
For a fixed realization, the displayed empirical process is
$\int f\,d\sigma_ N$. The Lipschitz class is invariant under $f\mapsto-f$,
so the absolute value does not change the supremum. The last equality is
the equal-mass Kantorovich--Rubinstein duality. 
\end{proof}

\subsection{Comparison with the Gaussian grid}

\begin{definition}[Balanced Gaussian grid]\label{def:gaussian-grid}
Partition $[0,1]^2$ into a $K\times K$ grid, let $c_ z$ be the cell
centres, let $G_ z$ be independent standard Gaussians, and put
\begin{equation}
\bar G=K^{-2}\sum_ zG_ z,\qquad
\Gamma_ K=\mathbb E\left\|
\sum_ z(G_ z-\bar G)\delta_ {c_ z}\right\|_ {\rm KR}.
\tag{4}
\end{equation}

Subtraction of $\bar G$ makes the total mass zero.

\end{definition}

\begin{lemma}[Quantitative Gaussian approximation]\label{lem:gaussianization-dyadic}
Uniformly for $K^2\leq N$,
\begin{equation}
\left|A_ N-\sqrt{\frac N{K^2}}\,\Gamma_ K\right|
\leq C\left(\sqrt N+\frac NK+K\sqrt{\log(2K)}\right).
\tag{5}
\end{equation}
Consequently, it is enough to prove the existence of
\begin{equation}
\lim_ {J\to\infty}
\frac{\Gamma_ {2^J}}{2^J\sqrt{J\log2}}.
\tag{6}
\end{equation}
If this dyadic limit is $\gamma$, then
\begin{equation}
\lim_ {N\to\infty}\frac{A_ N}{\sqrt{N\log N}}
=\frac\gamma{\sqrt2}.
\tag{7}
\end{equation}

\end{lemma}

\begin{proof}
Poissonization and snapping produce independent cell counts at costs
$O(\sqrt N)$ and $O(N/K)$.  After coupling these counts to Gaussians,
we construct a grid flow for the balanced error; its energy controls the
third term in (5).  The final choice of $K$ makes all three errors
negligible on the scale $\sqrt{N\log N}$.

\smallskip\noindent
\underline{\textit{Step 1.}} Poissonize and snap to the grid.

Write $\mu_ n=\sum_ {i=1}^n(\delta_ {X_ i}-\lambda)$.  Couple
$M\sim\operatorname {Pois}(N)$ to the same infinite sample sequence.
Since $\|\delta_ x-\lambda\|_ {\rm KR}\leq\operatorname {diam}Q_ 1$,
\begin{equation}
\mathbb E\|\mu_ N-\mu_ M\|_ {\rm KR}
\leq C\mathbb E|M-N|\leq C\sqrt N.
\tag{5a}
\end{equation}
Let $\lambda_ K=K^{-2}\sum_ z\delta_ {c_ z}$.  Replacing each
$\delta_ {X_ i}-\lambda$ by
$\delta_ {c(X_ i)}-\lambda_ K$ costs at most $C/K$.  Thus the expected
snapping error in the Poissonized process is $CN/K$.

\smallskip\noindent
\underline{\textit{Step 2.}} Couple the Poisson counts to Gaussians.

The cell counts $P_ z$ are now independent
$\operatorname {Pois}(q)$ variables, $q=N/K^2$, while the balanced
snapped measure is
\begin{equation}
\sum_ z(P_ z-\bar P)\delta_ {c_ z},\qquad
\bar P=K^{-2}\sum_ zP_ z.
\tag{5b}
\end{equation}
Put $Y_ z=P_ z-q$. Couple the $Y_ z$'s independently to
$\sqrt q\,G_ z$
with
\begin{equation}
\mathbb E|P_ z-q-\sqrt q\,G_ z|^2\leq C.
\tag{8}
\end{equation}
We use the following normal-approximation estimate to obtain (8).

\textbf{Normal-approximation result.} If $V_ n=\sum_ i\xi_ i$, the $\xi_ i$
are independent and centred, and $\mathbb EV_ n^2=1$, then for every
$p>1$
\begin{equation*}
W_ p(\mathcal L(V_ n),N(0,1))
\leq C_ p\left(\sum_ i\mathbb E|\xi_ i|^{p+2}\right)^{1/p}.
\end{equation*}
Here $C_ p$ depends only continuously on $p$.

This is Lemma 3.4 of Fang~\cite{Fang2019}, which restates Bobkov's Theorem 1.1. To check the hypotheses without a limiting argument, let
$r=\lceil q\rceil$, $\theta=q/r\in[1/2,1]$, and use infinite divisibility to
write $P_ q$ as a sum of $r$ independent $\mathrm {Pois}(\theta)$
variables.  After normalization to variance one, the sum of fourth
moments appearing in the $p=2$ estimate is
$(1+3\theta)/q\leq4/q$.  Hence
$W_ 2((P_ q-q)/\sqrt q,N(0,1))\leq C/\sqrt q$, which is exactly
(8), uniformly for $q\geq1$.

Let $e_ z=Y_ z-\sqrt qG_ z$ and
$\widetilde e=e-\bar e\mathbf1$.  The couplings are independent and
centred, so
$\operatorname {Cov}(e)\leq CI$ and
$\operatorname {Cov}(\widetilde e)\leq C\Pi_ 0$, where $\Pi_ 0$ is
orthogonal projection off the constants.  Notice that projection is
exactly what turns $Y_ z$ into $P_ z-\bar P$ and $G_ z$ into
$G_ z-\bar G$; it introduces no unbalanced measure.

\smallskip\noindent
\underline{\textit{Step 3.}} Transport the balanced coupling error.

The errors in (8) have uniformly bounded variances.  Summing their
absolute values over $K^2$ cells would discard cancellation.  We instead
use the discrete inverse Laplacian to construct a flow whose divergence
is exactly the balanced error.

Give
the $K\times K$ nearest-neighbor grid unit edge weights and let
$B:\mathbb R^E\to\mathbb R^{K^2}$ be its signed incidence matrix,
$L=BB^T$.  The flow
\begin{equation*}
m=B^TL^+\widetilde e
\end{equation*}
satisfies $Bm=\widetilde e$.  Sending $|m_ e|$ units of mass along
the corresponding Euclidean edge of length $1/K$ gives
\begin{equation}
\left\|\sum_ z\widetilde e_ z\delta_ {c_ z}\right\|_ {\rm KR}
\leq K^{-1}\sum_ {e\in E}|m_ e|.
\tag{5c}
\end{equation}
Moreover
\begin{equation*}
\mathbb E\|m\|_ 2^2
=\operatorname {Tr}\!\left(B^TL^+\operatorname {Cov}(\widetilde e)
L^+B\right)
\leq C\operatorname {Tr}L^+.
\end{equation*}
For the path-grid Neumann Laplacian its nonzero eigenvalues are
\begin{equation*}
\lambda_{r,s}=4\sin^2\!\left(\frac{\pi r}{2K}\right)+4\sin^2\!\left(\frac{\pi s}{2K}\right),
\end{equation*}
with the evident one-axis variants.  Therefore
\begin{equation}
\operatorname {Tr}L^+
\leq CK^2\sum_ {0\leq r,s<K,(r,s)\ne(0,0)}(r^2+s^2)^{-1}
\leq CK^2\log(2K).
\tag{5d}
\end{equation}
Since $|E|\leq2K^2$, Cauchy--Schwarz in (5c) yields expected KR error
at most $CK\sqrt{\log(2K)}$.  Together with (5a) and the snapping
bound this proves (5).

\smallskip\noindent
\underline{\textit{Step 4.}} Choose the dyadic grid scale and return to $A_N$.

Fix $0<a<1/2$. For every $N$, choose a dyadic integer $K_ N$ such
that
\begin{equation*}
K_ N\leq\frac{\sqrt N}{(\log N)^a}<2K_ N.
\end{equation*}
All three errors in (5) are $o(\sqrt{N\log N})$, while
$\log K_ N/\log N\to1/2$.  Substitution of (6) into (5) therefore
proves (7), completing the reduction from the empirical problem to the
dyadic Gaussian limit.

\end{proof}

\subsection{From Gaussian atoms to a cellwise density}

The next comparison puts the Gaussian field on a square with unit cells,
where its dyadic decomposition is convenient.  The resulting error is of
lower order than the target cost $K^2\sqrt{\log K}$.

\begin{lemma}[Spreading the Gaussian masses over cells]\label{lem:atom-to-cell}
Let $\widehat\Gamma_ K$ be the Gaussian KR expectation after dilating the
unit square and the grid by $K$. Thus
\begin{equation*}
\widehat\Gamma_ K=K\Gamma_ K.
\end{equation*}
Let $C_ z=z+[0,1)^2\subset Q_ K$ and set
\begin{equation*}
\rho_ K(x)=\sum_ z(G_ z-\bar G)\mathbf1_ {C_ z}(x).
\end{equation*}
Then
\begin{equation}
\left|\widehat\Gamma_ K-
\mathbb E\|\rho_ K\|_ {\rm KR,Q_ K}\right|\leq CK^2.
\tag{9}
\end{equation}

\end{lemma}

\begin{proof}
Dilation multiplies every transport distance by $K$, proving the first
identity. For each cell, move its signed mass between the centre and the
uniform measure on that cell. This costs at most
$C|G_ z-\bar G|$. Summing and taking expectations proves (9).

\end{proof}

\section{Dyadic decomposition and scale grouping}\label{sec:dyadic}

We group the Gaussian density into independent bands of dyadic scales.
Each band has the same law after its outer blocks are rescaled to unit
squares.  To separate consecutive retained bands, we will omit some Haar
levels; the estimate below shows that the resulting cost grows only as
the square root of their number.  Its proof constructs local fluxes whose
pointwise second moments also control the later cell comparisons.

\begin{lemma}[Independent Haar bands and the cost of omitted levels]\label{lem:dyadic-bands}
Assume $K=2^J$. On cellwise-constant functions in $Q_ K$, let $P_ j$
be conditional expectation on dyadic squares of side $2^j$. Then
\begin{equation}
\rho_ K=\sum_ {j=0}^{J-1}\xi_ j,
\qquad \xi_ j=(P_ j-P_ {j+1})G,
\tag{10}
\end{equation}
and the $\xi_ j$ are independent Gaussian fields.

If $h\geq1$, then a group
$\eta_ {r,h}=(P_ r-P_ {r+h})G$ is balanced on every dyadic square of side
$2^{r+h}$, and different such outer squares carry independent copies.
After rescaling an outer square to the unit square, its source is
\begin{equation}
\eta_ h=n\sum_ {C\in\mathcal P_ n}(Z_ C-\bar Z)\mathbf1_ C,
\qquad n=2^h,
\tag{11}
\end{equation}
where $\mathcal P_ n$ is the $n\times n$ partition into cells of side
$1/n$, the $Z_ C$ are i.i.d. $N(0,1)$, and
$\bar Z=n^{-2}\sum_ CZ_ C$. In particular, the rescaled law depends on
$h$, but not on $r$.

We call a square of side $2^{r+h}$ on which $\eta_{r,h}$ is balanced an
\emph{outer block} of the band.  After rescaling, an outer block has side
one and contains $2^h\times2^h$ fine cells.

If $D$ dyadic levels are omitted, their sum $\zeta_ D$ satisfies
\begin{equation}
\mathbb E\|\zeta_ D\|_ {\rm KR,Q_ K}\leq CK^2\sqrt D.
\tag{12}
\end{equation}

\end{lemma}

\begin{proof}
By (1), it suffices to construct a flux for the omitted levels with
expected $L^1$ norm at most $CK^2\sqrt D$.  We construct a zero-normal
flux for each Haar detail and use independence to add its second moment
across levels.  First we identify the law of each band.

\smallskip\noindent
\underline{\textit{Step 1.}} Decompose the field into independent dyadic bands.

The spaces $\operatorname {ran}(P_ j-P_ {j+1})$ are mutually orthogonal.
Orthogonal projections of a Gaussian vector are independent, proving
(10). The balance and independence assertions follow directly from the
definition of conditional expectation. Under $x=2^{r+h}y$, fluxes are
unchanged in dimension two and source densities acquire the factor
$2^{r+h}$; this gives exactly (11).

\smallskip\noindent
\underline{\textit{Step 2.}} Construct a local right inverse on one Haar parent.

The local flux must solve the divergence equation and have matching
normal components at every child interface; otherwise pasting would
create an additional source there.  Work on one level-$j$ parent
$B=x_ 0+[0,2\ell)^2$, $\ell=2^j$.  Denote its four child averages by
$A_ {\sigma\tau}$, where $\sigma=1$ means left, $\tau=1$ means lower,
and $\sigma,\tau\in\{-1,1\}$.  They are independent
$N(0,\ell^{-2})$.  Put
\begin{equation*}
\alpha=\frac14\sum_ {\sigma,\tau}\sigma A_ {\sigma\tau},\qquad
\beta =\frac14\sum_ {\sigma,\tau}\tau A_ {\sigma\tau},\qquad
\gamma=\frac14\sum_ {\sigma,\tau}\sigma\tau A_ {\sigma\tau}.
\end{equation*}
Thus $\alpha,\beta,\gamma$ are independent and each has variance
$(4\ell^2)^{-1}$.  Let $s_ \ell=1$ on $[0,\ell)$ and $-1$ on
$[\ell,2\ell)$, and let
\begin{equation*}
\tau_ \ell(t)=\begin{cases}t,&0\leq t\leq\ell,\\
2\ell-t,&\ell<t\leq2\ell,
\end{cases}
\qquad \tau_ \ell'=s_ \ell\quad\hbox{a.e.}
\end{equation*}
The Haar detail on $B$ is
\begin{equation*}
\xi_ j=\alpha s_ \ell(x_ 1-x_ {0,1})+\beta s_ \ell(x_ 2-x_ {0,2})
+\gamma s_ \ell(x_ 1-x_ {0,1})s_ \ell(x_ 2-x_ {0,2}),
\end{equation*}
and define
\begin{equation}
\begin{split}
(R_ j\xi_ j)_ 1(x)&=-\tau_ \ell(x_ 1-x_ {0,1})
\{\alpha+\gamma s_ \ell(x_ 2-x_ {0,2})\},\\
(R_ j\xi_ j)_ 2(x)&=-\beta\tau_ \ell(x_ 2-x_ {0,2}).
\end{split}
\tag{12a}
\end{equation}
Direct differentiation gives $-\operatorname {div}R_ j\xi_ j=\xi_ j$.
The first component vanishes on the two vertical parent faces and the
second on the two horizontal parent faces, so the outer normal trace is
zero.  At the internal vertical child interface the normal first
component is continuous; at the internal horizontal interface the normal
second component is continuous, while only the tangential first component
jumps.  Hence there is no interface measure in the divergence.

\smallskip\noindent
\underline{\textit{Step 3.}} Paste the local fields and retain a pointwise bound.

Paste (12a) over all level-$j$ parents.  Across a common parent face both
normal traces are zero, so the pasted field lies in
$H(\operatorname {div})$ and again creates no interface source.  Choose
the representative zero on the union of dyadic interfaces.  Since
$0\leq\tau_ \ell\leq\ell$, at every point
\begin{equation}
\mathbb E|R_ j\xi_ j(x)|^2\leq \frac34,\qquad
\mathbb E\int_ {\text{parent}}|R_ j\xi_ j|^2\leq C\ell^2.
\tag{12b}
\end{equation}
The Gaussian Haar projections at different levels are orthogonal and
therefore independent; $R_ j$ is deterministic and linear.  Consequently,
for every set $I$ of levels,
\begin{equation}
\sup_ x\mathbb E\left|\sum_ {j\in I}R_ j\xi_ j(x)\right|^2
=\sup_ x\sum_ {j\in I}\mathbb E|R_ j\xi_ j(x)|^2
\leq C|I|.
\tag{12c}
\end{equation}
The pointwise estimate (12c), including the chosen representatives at
interfaces, is used in (16), (17), and (22).  The construction is unchanged
by dilation of an outer block: fluxes are scale invariant in dimension two.

\smallskip\noindent
\underline{\textit{Step 4.}} Sum the omitted levels and conclude the KR estimate.

For the $D$ omitted levels, the same independence gives
\begin{equation*}
\mathbb E\int_ {Q_ K}\left|\sum_ {j\ {\rm omitted}}R_ j\xi_ j\right|^2
=\sum_ {j\ {\rm omitted}}\mathbb E\int_ {Q_ K}|R_ j\xi_ j|^2
\leq CDK^2.
\end{equation*}
Jensen followed by Cauchy--Schwarz in space bounds the expected
$L^1$ norm by $CK^2\sqrt D$.  Formula (1) now gives (12), completing
the proof.

\end{proof}

\section{The nonlinear cell operator}\label{sec:cell-operator}

After the finer scales have been treated, their contribution is encoded
by a convex cost profile $\Psi(p)$ for a base flux $p$.  The operator
defined below incorporates one more Haar band and returns the updated
profile.  We define it through the dual potential: first prescribe its
values on all four cell faces, optimize the interior for each fine-source
realization, and take expectation; only then optimize the boundary trace.
This order keeps the trace deterministic with respect to the fine source.
The next subsection estimates the gain from allowing that trace to depend
on the fine source.

\subsection{Definition and basic properties}

\begin{definition}[Cell operator]\label{def:cell-operator}
Let $\mathscr P_ M$ be the class of convex functions
$\Psi:\mathbb R^2\to\mathbb R$ satisfying
\begin{equation}
|p|\leq\Psi(p)\leq |p|+M.
\tag{13}
\end{equation}
We call the members of $\mathscr P_M$ \emph{linear-growth profiles}.
Every member of $\mathscr P_ M$ is $1$-Lipschitz, and
\begin{equation}
-M\leq\Psi^*(a)\leq0\quad (|a|\leq1),\qquad
\Psi^*(a)=+\infty\quad (|a|>1).
\tag{14}
\end{equation}

For an integer $L\geq1$, tile $Q_ L$ by unit squares (the normalized outer
blocks) and put an independent copy of (11) in each square; call the resulting
source $\eta_ h^{(L)}$.  To retain the boundary data needed for pasting
cell interiors, let
$\mathcal T_ L$ be the restrictions to all four faces of $\partial Q_ L$
of $1$-Lipschitz functions on $Q_ L$, modulo constants and represented
by the normalization $c(0)=0$.  Since $Q_ L$ is convex, McShane
extension shows that this is exactly the compact set of anchored
$1$-Lipschitz boundary functions.

The slope $a$ will record the mean gradient of the dual potential.
For $|a|\leq1$, let $\mathcal A_ L(a)$ be the functions
$F\in W^{1,\infty}_ {\rm loc}(\mathbb R^2)$ satisfying
\begin{equation}
|\nabla F|\leq1,\qquad
F(x+Le_ i)=F(x)+La_ i\quad(i=1,2),\qquad F(0)=0.
\tag{14a}
\end{equation}
Let $\mathcal T_ L^{\rm per}(a)\subset\mathcal T_ L$ be the restrictions
of these functions to all four faces of $\partial Q_ L$.  A
\emph{complete affine-periodic trace} records the actual height on every
face, the opposite-face increment relations $F(x+Le_i)-F(x)=La_i$,
and the corner compatibility.  These are the data that allow one cell
extension to be repeated over the plane.  For every $c\in\mathcal T_ L$,
let
\begin{equation}
K_ {h,L,\Psi}(c)=\mathbb E\sup_ {\substack{F|_ {\partial Q_ L}=c\\
|\nabla F|\leq1}}
\left\{\int_ {Q_ L}F\eta_ h^{(L)}
-\int_ {Q_ L}\Psi^*(\nabla F)\right\}.
\tag{15}
\end{equation}
The source has zero total mass, so adding a constant to $c$ leaves this
value unchanged.  The best expected cell value at slope $a$ and the
updated cost at base flux $p$ are, respectively,
\begin{equation}
\mathcal R_ {h,L,\Psi}(a)=L^{-2}
\max_ {c\in\mathcal T_ L^{\rm per}(a)}K_ {h,L,\Psi}(c),\qquad
(\mathcal H_ {h,L}\Psi)(p)
=\sup_ {|a|\leq1}\{p\cdot a+\mathcal R_ {h,L,\Psi}(a)\}.
\tag{16}
\end{equation}
When $h,L,\Psi$ are fixed, we abbreviate
$\mathcal R=\mathcal R_{h,L,\Psi}$.

\smallskip\noindent\emph{Extension to a compact slope domain.} Let
$\Psi$ be proper closed convex and suppose
$K_ \Psi=\operatorname {dom}\Psi^*\Subset B_ 1$ is compact and
$\Psi^*$ is finite and continuous on $K_ \Psi$.  In (15), retain the
same ambient trace space but interpret the value as $-\infty$ unless
the trace has an extension with
$\nabla F(x)\in K_ \Psi$ almost everywhere; equivalently the extended
integrand $\Psi^*=+\infty$ off $K_ \Psi$ enforces this condition.
In (16), set $\mathcal R(a)=-\infty$ for $a\notin K_ \Psi$.  These conventions
define $\mathcal H_ {h,L}\Psi=(-\mathcal R)^*$.  This extension is used only for the strict-slope lower barrier in
Lemma~\ref{lem:viscosity-limit}.

\end{definition}

\begin{proposition}[Properties of the cell operator]\label{prop:cell-operator-properties}
On the linear-growth class in Definition~\ref{def:cell-operator}, the map $\mathcal H_ {h,L}$ is order preserving,
commutes with addition of constants, is a contraction in the uniform
norm, and maps convex $1$-Lipschitz profiles to convex
$1$-Lipschitz profiles.  Moreover,
\begin{equation}
\Psi\leq\mathcal H_ {h,L}\Psi\leq\Psi+C\sqrt h.
\tag{16a}
\end{equation}
The bound is uniform in $L$.  If $\Psi$ is invariant under the dihedral group
$D_ 4$ of the square, so is $\mathcal H_ {h,L}\Psi$.
On the extension to a compact slope domain in Definition~\ref{def:cell-operator},
the map is also order preserving and constant-equivariant.
\end{proposition}

\begin{proof}
We first establish the variational properties on $\mathscr P_M$ and
bound the update by the expected size of a Haar flux.  We then verify
that the same definition remains well posed on a compact slope domain.

\smallskip\noindent
\underline{\textit{Step 1.}} Establish the variational properties on $\mathscr P_M$.

The bounds in (14) follow from (13) and convex conjugacy. A convex
function whose recession function is $|\cdot|$ has all subgradients in
the closed unit ball and is therefore $1$-Lipschitz.

The space $\mathcal T_ L$, and the closed set of pairs $(a,c)$, with
$|a|\leq1$ and $c\in\mathcal T_ L^{\rm per}(a)$, are compact in the
uniform topology by Arzela--Ascoli.  Repeating a trace and any one of its
cell extensions produces a continuous piecewise $W^{1,\infty}$ function
on the plane; its weak gradient has norm at most one, so it is globally
$1$-Lipschitz.  Thus the boundary cocycle conditions are also sufficient,
not only necessary.  For fixed trace, the admissible interior potentials are
also compact after normalization.  The linear source term is uniformly
continuous and $F\mapsto-\int\Psi^*(\nabla F)$ is upper
semicontinuous under weak-* convergence in $W^{1,\infty}$.  Hence the
two maxima in (15)--(16) are attained.  Mixing traces and
near-optimizing interiors shows that $\mathcal R$ is concave and upper
semicontinuous. Thus (16) is convex and $1$-Lipschitz and
$(\mathcal H_ {h,L}\Psi)^*=-\mathcal R$ on the closed unit ball.  If
$\Psi\leq\Phi$, then $\Psi^*\geq\Phi^*$, which proves order
preservation; the constant rule and contraction follow in the same way.

\smallskip\noindent
\underline{\textit{Step 2.}} Bound the update and preserve the symmetries.

The affine trace $F(x)=a\cdot x$ gives
$\mathcal R(a)\geq-\Psi^*(a)$, so biconjugacy gives the lower bound in
(16a).  For the upper bound, we convert the source pairing into a perturbation
of the base flux.  Paste on the unit outer blocks the linear zero-normal
flux $Q$ from (12a), so that
$-\operatorname {div}Q=\eta_ h^{(L)}$,
$\mathbb E Q(x)=0$, and
$\sup_ x\mathbb E|Q(x)|^2\leq Ch$.  For every affine-cocycle potential
$F$ of slope $a$, integration by parts and
$\fint\nabla F=a$ give
\begin{equation}
p\cdot a+L^{-2}\left(\int F\eta_ h^{(L)}-
\int\Psi^*(\nabla F)\right)
=\fint\bigl[(p+Q)\cdot\nabla F-\Psi^*(\nabla F)\bigr]
\leq\fint\Psi(p+Q).
\tag{16b}
\end{equation}

Taking the fixed-trace maximum and expectation in (16b), then using that
$\Psi$ is $1$-Lipschitz, yields
\begin{equation}
(\mathcal H_ {h,L}\Psi)(p)
\leq\mathbb E\fint\Psi(p+Q)\leq\Psi(p)+C\sqrt h.
\tag{16c}
\end{equation}
This proves (16a) directly from the dual cell problem.
The same argument, without the $p$-term, also gives
\begin{equation}
0\leq L^{-2}K_ {h,L,\Psi}(c)\leq C(M+\sqrt h),
\qquad 0\leq \mathcal R_ {h,L,\Psi}(a)\leq C(M+\sqrt h).
\tag{16d}
\end{equation}
For the lower bound, use any deterministic extension of $c$: its
expected source pairing is zero and $-\Psi^*\geq0$.  The upper bound
follows from integration by parts against the pasted local flux $Q$
and Fenchel's inequality.  These estimates also imply the integrable
domination needed when traces converge.

Applying a coordinate reflection or permutation simultaneously to the
source, flux, and trace preserves their laws and the variational value;
this proves the $D_ 4$-invariance assertion.

\smallskip\noindent
\underline{\textit{Step 3.}} Treat a compact slope domain.

The additional issue is to retain the constraint
$\nabla F\in K_\Psi$ when traces and interior potentials converge.
Write $g=\Psi^*$, extended by
$+\infty$ off $K_ \Psi$.  The compatible traces, and the set of
compatible affine-cocycle pairs $(a,c)$, are closed.  Indeed, if
$F_ n$ are normalized extensions with
$\nabla F_ n\in K_ \Psi$, uniform convergence and weak-* compactness in
$W^{1,\infty}$ give an extension $F$.  For every affine functional
$\ell$ nonpositive on $K_ \Psi$, weak-* passage in
$\int_ E\ell(\nabla F_ n)$ gives
$\int_ E\ell(\nabla F)\leq0$ for every measurable $E$;
finite-dimensional separation therefore gives
$\nabla F\in K_ \Psi$ almost everywhere.  The trace and cocycle
identities pass to the limit.  Thus both compatible sets are compact.

The functional $\int g(\nabla F)$ is weak-* lower semicontinuous, so
the payoff in (15) is upper semicontinuous and the fixed-trace maximum
is attained.  For a fixed source its value is upper semicontinuous in
the compatible trace: take maximizing extensions along a convergent
sequence of traces and use the preceding compactness.  The maximum is
measurable in the finite Gaussian coefficient vector by the measurable
maximum theorem.  Moreover, after anchoring one value,
\begin{equation}
\left|\max_ F\left\{\int F\eta_ h^{(L)}-\int g(\nabla F)\right\}\right|
\leq C_ {K_ \Psi,L,g}\bigl(1+\|\eta_ h^{(L)}\|_ {L^1(Q_ L)}\bigr),
\tag{16s}
\end{equation}
uniformly over compatible traces.  The right side is integrable.
Reverse Fatou therefore shows that the expected fixed-trace value is
upper semicontinuous.  Compactness of the affine-cocycle pairs now
proves attainment and upper semicontinuity of the outer maximum
defining $\mathcal R$.

Convex combinations of compatible potentials remain compatible, and
the payoff is concave in the potential.  Mixing affine-cocycle
potentials proves that $\mathcal R$ is concave.  Every compatible potential of
slope $a$ satisfies
\begin{equation*}
a=\fint_ {Q_ L}\nabla F\in K_ \Psi,
\end{equation*}
whereas $F(x)=a\cdot x$ is compatible for every $a\in K_ \Psi$.
Consequently $\mathcal R$ is finite precisely on $K_ \Psi$, is
upper semicontinuous there, and equals $-\infty$ off it.  Hence
$-\mathcal R$ is a proper closed convex extended-real function and
$\mathcal H_ {h,L}\Psi=(-\mathcal R)^*$, as asserted.

Incompatible traces have value $-\infty$.  Thus, if
$\Psi\leq\Phi$, extended-real conjugacy gives
$\Psi^*\geq\Phi^*$ pointwise even when their compact domains differ;
the payoff inequality may be maximized over the common ambient trace
space and proves successively the order of the fixed-trace values,
$\mathcal R$, and $\mathcal H$.  Finally
$(\Psi+C)^*=\Psi^*-C$, which proves constant equivariance.

\end{proof}

\subsection{A finite approximation of common boundary traces}

A global optimizing potential can use the fine source to choose its trace
on every cell boundary.  To compare this value with the cell operator,
we approximate all of those boundary values simultaneously by a finite
family of compatible traces.  For each fixed trace the cell interiors
optimize separately, and Gaussian concentration bounds the gain from
choosing the trace after seeing the source.  The lemma is stated for a
constant base flux; the proof also records the conditional estimates
needed when that constant is chosen from an independent coarse field.

\begin{lemma}[Common-trace comparison]\label{lem:finite-trace-net}
Fix an $h$-level source and a convex profile $\Psi\in\mathscr P_ M$.
Let $\Omega=Q_ R$, with $R/L$ an integer, and let $\Sigma_ L$ be
the union of all boundaries of its $L$-squares.  A complete common
trace means the restriction of one global $1$-Lipschitz potential to
all of $\Sigma_ L$, modulo one additive constant; denote this compact
trace space by $\mathfrak T_ {\Omega,L}$.  Let $\eta$ consist of independent
copies of $\eta_h^{(L)}$ in these cells.  Fix a deterministic constant
base vector $p\in\mathbb R^2$ and define
\begin{equation}
\mathcal V_ {\Omega}(p,\eta)=
\sup_ {\substack{F\in W^{1,\infty}(\Omega),\ |\nabla F|\leq1\\F(0)=0}}
\left\{\int_ \Omega F\eta+\int_ \Omega p\cdot\nabla F
-\int_ \Omega\Psi^*(\nabla F)\right\}.
\tag{17a}
\end{equation}

For $c\in\mathfrak T_ {\Omega,L}$, let
$\mathcal V_ c(p,\eta)$ be (17a) with $F|_ {\Sigma_ L}=c$.
If $D$ is an $L$-cell, write $c_ D$ for
the restricted trace and
\begin{equation}
m_ D(c)=\int_ D\nabla F
=\int_ {\partial D}c_ D\,n\,d\mathcal H^1,
\tag{17b}
\end{equation}
which depends only on the trace.  After translation to the reference
cell we denote the same quantity by $m_ L(c_ D)$.  For
$0<\varepsilon<1/10$ there is
a finite net $\mathcal N\subset\mathfrak T_ {\Omega,L}$ such that
\begin{equation}
\begin{split}
\mathbb E\mathcal V_ \Omega(p,\eta)
\leq{}&\max_ {c\in\mathcal N}
\sum_ D\{p\cdot m_ D(c)+K_ {h,L,\Psi}(c_ D)\}\\
&+|\Omega|\left[
C(M+|p|+\sqrt h)\left(\varepsilon+\frac1{\varepsilon L}\right)
+C\sqrt{\frac{\log(C/\varepsilon)}L}\right].                   
\end{split}
\tag{17}
\end{equation}
Conversely, for every deterministic common trace $c$,
\begin{equation}
\mathbb E\mathcal V_ \Omega(p,\eta)
\geq\sum_ D\{p\cdot m_ D(c)+K_ {h,L,\Psi}(c_ D)\}.
\tag{17c}
\end{equation}
In Lemma~\ref{lem:two-scale-scalarization}, the coarse field is first
frozen to a constant $p_D$ on each block before this estimate is applied.

For a fixed $L$-cell and an affine cocycle $a$, let $K(c)$ be the
expected optimized fine dual value with complete boundary trace $c$.
Then $K$ is concave in $c$. If
\begin{equation}
\mathcal R(a)=L^{-2}\max\{K(c):c\text{ has affine-periodic slope }a\},
\qquad |a|\leq1,
\tag{18}
\end{equation}
then $\mathcal R$ is concave and upper semicontinuous and
\begin{equation}
H(p)=\sup_ {|a|\leq1}\{p\cdot a+\mathcal R(a)\},\qquad H^*=-\mathcal R
\ \hbox{ on }\overline B_ 1.
\tag{19}
\end{equation}
With the notation of Definition~\ref{def:cell-operator},
$H=\mathcal H_ {h,L}\Psi$ exactly.  The cell size remains finite throughout this comparison.

\end{lemma}

\begin{proof}
We prove (17) by approximating the random trace and then comparing the
maximum over that finite family with its expected fixed-trace values.
To make its later conditional use explicit, we carry out the construction
with an independent coarse field.

On a coarse probability space let $P$ be a strongly
$L^1(\Omega;\mathbb R^2)$-measurable field with
$\mathbb E\int_ \Omega|P|<\infty$.  Independently, on a product fine
probability space, put independent copies $\eta_ D$ of the active
fine source in the $L$-cells and set $\eta=\sum_ D\eta_ D$.  With
\begin{equation*}
\mathscr A_ \Omega=\{F\in C(\overline\Omega):F(0)=0,
\operatorname {Lip}(F)\leq1\},
\end{equation*}
set
\begin{equation}
\mathcal V_ \Omega(P,\eta)=
\max_ {F\in\mathscr A_ \Omega}
\left\{\langle\eta,F\rangle+
\int_ \Omega P\cdot\nabla F-
\int_ \Omega\Psi^*(\nabla F)\right\},\qquad
\mathbf D_ {R,L}=\mathbb E_ P\mathbb E_ \eta
\mathcal V_ \Omega(P,\eta).
\tag{17a'}
\end{equation}
The maximum is the free-outer-boundary dual problem; no periodic outer
trace is imposed.  The displayed integrability assumption and the
canonical fine flux bound make $\mathbf D_ {R,L}$ finite.

Let $\mathcal V_c(P,\eta)$ denote (17a') with
$F|_{\Sigma_L}=c$.  For $P\equiv p$, these definitions agree with
those in the statement.

\smallskip\noindent
\underline{\textit{Step 1.}} Construct a finite family of compatible traces.

By (2), $\mathcal V_ \Omega$ is exactly the free-boundary dual
supremum in (17a').
Writing the fine balanced source as $-\operatorname {div}Q$, its
objective is
\begin{equation}
\sup_ F\int\bigl[(P+Q)\cdot\nabla F-\Psi^*(\nabla F)\bigr],
\qquad |\nabla F|\leq1.
\tag{20}
\end{equation}
Rounding and interpolating each face separately would not ensure agreement
at junctions or compatibility around cycles.  We therefore round values
on the whole skeleton and use one McShane extension.  The preliminary
contraction by $\theta=1-4\varepsilon$ supplies the slack for rounding.

Normalize $F$ at one skeleton vertex.  Scale it by
$\theta=1-4\varepsilon$, sample $\Sigma_ L$ at unit spacing (including
every junction), and round the sampled values to a lattice of mesh
$\varepsilon/8$.  Distinct sample vertices have Euclidean distance at
least one, and therefore their rounded values $q(v)$ satisfy
\begin{equation*}
|q(v)-q(w)|\leq\theta|v-w|+\varepsilon/4
\leq(1-2\varepsilon)|v-w|.
\end{equation*}
Define the whole skeleton trace, rather than interpolating face by face, by
\begin{equation}
c_ q(x)=\inf_ v\{q(v)+(1-2\varepsilon)|x-v|\},
\qquad x\in\Sigma_ L.
\tag{20a}
\end{equation}
This McShane formula agrees with the node data and is globally
$(1-2\varepsilon)$-Lipschitz in the Euclidean metric, so every cycle,
junction, and seam compatibility is automatic.  Comparison with the
nearest node on the same unit skeleton edge gives
$\|c_ q-\theta F\|_ {L^\infty(\Sigma_ L)}\leq1$.
These traces form the net $\mathcal N$.
The skeleton graph has at most $C|\Omega|/L$ unit edges.  Once one
anchor is fixed, each new vertex has at most $C/\varepsilon$ possible
increments (cycle constraints only reduce this number).  Hence
\begin{equation}
\log|\mathcal N|
\leq C\frac{|\Omega|}{L}\log(C/\varepsilon).
\tag{21}
\end{equation}
\smallskip\noindent
\underline{\textit{Step 2.}} Match the optimizer to the common trace.

Preserve the entire scaled optimizer on the core at distance
$r=C/\varepsilon$ from the skeleton. For a core
point $x$ and a skeleton point $y$,
\begin{equation*}
|F_ {\rm core}(x)-F_ {\rm skel}(y)|
\leq C+(1-4\varepsilon)|x-y|\leq(1-\varepsilon)|x-y|.
\end{equation*}
Thus the union of all core data and the common skeleton trace has a
McShane $1$-Lipschitz extension.  Fine sources in the intervening
corridors are compared using their canonical unit-block zero-normal
fluxes.  Let $\mathcal C_ {\Omega,L,\varepsilon}$ be the resulting
union of whole unit outer blocks.  It has deterministic relative area at most
$C/(\varepsilon L)$, and, because every unit-outer-block source is balanced,
\begin{equation}
\mathbb E_ \eta\int_ {\mathcal C_ {\Omega,L,\varepsilon}}|Q|
\leq C|\Omega|\frac{\sqrt h}{\varepsilon L}.
\tag{21a}
\end{equation}
On the core, concavity of $-\Psi^*$, together with
$-M\leq\Psi^*\leq0$, bounds the loss by
$C\varepsilon(M+|P|+|Q|)$.  On the corridor, both old and new
integrands have magnitude at most $|P|+|Q|+M$.  Consequently this modification
gives, for every realized coarse field $P$,
\begin{equation}
\begin{split}
0\leq{}&\mathbb E_ \eta[\mathcal V_ \Omega(P,\eta)
-\max_ {c\in\mathcal N}\mathcal V_ c(P,\eta)\mid P]\\
\leq{}&C|\Omega|(M+\sqrt h)
\left(\varepsilon+\frac1{\varepsilon L}\right)
+C\varepsilon\int_ \Omega|P|
+C\int_ {\mathcal C_ {\Omega,L,\varepsilon}}|P|.          
\end{split}
\tag{21a'}
\end{equation}
The last two terms retain the spatial dependence of the realized field
$P$.  If
$\sup_ x\mathbb E|P(x)|\leq C\sqrt T$, only now averaging in $P$
and using the deterministic corridor fraction turns the last two terms
into $C|\Omega|\sqrt T(\varepsilon+1/(\varepsilon L))$.  For
$P\equiv p$, the same calculation gives the geometric error in (17).

\smallskip\noindent
\underline{\textit{Step 3.}} Decompose the fixed-trace problem cell by cell.

For a realized coarse field $P$, put
\begin{equation*}
\begin{aligned}
Y_ D(P,c_ D;\eta_ D)&=
\max_ {\substack{F_ D|_ {\partial D}=c_ D\\|\nabla F_ D|\leq1}}
\left\{\langle\eta_ D,F_ D\rangle+
\int_ DP\cdot\nabla F_ D-\int_ D\Psi^*(\nabla F_ D)\right\},\\
K_ {D,P}(c_ D)&=\mathbb E_ {\eta_ D}Y_ D(P,c_ D;\eta_ D).
\end{aligned}
\end{equation*}
Every family of cell interiors with one complete common trace pastes to
a global $1$-Lipschitz function: restrict the straight segment between
two points to its successive cell pieces and sum the cellwise Lipschitz
bounds.  Conversely every global potential restricts to such a family.
Thus, realization by realization,
\begin{equation}
\mathcal V_ \Omega(P,\eta)
=\max_ {c\in\mathfrak T_ {\Omega,L}}\sum_ DY_ D(P,c_ D;\eta_ D),
\qquad
\mathbb E_ \eta[\mathcal V_ c(P,\eta)\mid P]
=\sum_ DK_ {D,P}(c_ D).
\tag{21c}
\end{equation}
The second equality uses only the finite sum and independence of the
fine blocks.  In general $K_ {D,P}$ depends on the whole restriction
$P|_ D$, so it cannot be replaced by a function of $\fint_ DP$.
If $P=p_ D$ is constant on $D$, however, integration by parts gives
the exact specialization
\begin{equation}
K_ {D,P}(c_ D)=K_ {h,L,\Psi}(c_ D)+p_ D\cdot m_ D(c).
\tag{21d}
\end{equation}
In particular (21d) proves the fixed-trace identity used in (17), and
it also applies conditionally when the constants $p_ D$ are random but
measurable with respect to the independent coarse sigma-field.

\smallskip\noindent
\underline{\textit{Step 4.}} Control optimization over the finite trace family.

For a fixed skeleton trace, the optimized fine value in one $L$-cell
is $CL$-Lipschitz in its standard Gaussian coefficients.  Indeed, in
each unit outer block subtract the average of $F$; balance makes this
subtraction invisible, and the unit-block Poincare inequality gives
\begin{equation}
\sum_ {C\text{ fine}}
\left(n\int_ C(F-\fint_ {Q_ 1}F)\right)^2
\leq C\int_ {Q_ 1}|\nabla F|^2\leq C.
\tag{21b}
\end{equation}
Summing (21b) over the $L^2$ outer blocks gives Gaussian Lipschitz
radius $CL$.  Values in different $L$-cells depend on disjoint
Gaussian blocks.  Hence a fixed \emph{global} trace value has subgaussian
radius at most $C\sqrt{|\Omega|}$, independently of the realized
coarse field $P$.  The elementary
Gaussian estimate
\begin{equation*}
\mathbb E\max_ {\alpha\leq m}(Y_ \alpha-\mathbb EY_ \alpha)
\leq \sigma\sqrt{2\log m}
\end{equation*}
for $\sigma$-subgaussian variables, combined with (21), gives,
pointwise in $P$,
\begin{equation}
\mathbb E_ \eta\max_ {c\in\mathcal N}\mathcal V_ c(P,\eta)
\leq\max_ {c\in\mathcal N}\sum_ DK_ {D,P}(c_ D)
+C|\Omega|\sqrt{\frac{\log(C/\varepsilon)}L}.
\tag{21e}
\end{equation}
Specialize to $P\equiv p$, use (21d), and combine (21a') with (21e) to prove
(17).  Formula (17c) follows by restricting the original supremum to
any fixed common trace.  For random cellwise constants, (21d)--(21e) remain valid conditionally
on the coarse sigma-field.  We will take expectation in that field after
the deterministic comparison in the next section.

\smallskip\noindent
\underline{\textit{Step 5.}} Identify the deterministic one-cell operator.

If $c_ 1,c_ 2$ are two traces, convex combinations of their admissible
interiors are admissible, and the integrand in (20) is concave in
$\nabla F$. Hence $K$ is concave. On an $m\times m$ torus of
$L$-cells, average any deterministic trace pattern over all cell
translations on the universal cover, subtracting each translated value
at the origin before averaging.  Opposite sides telescope, so the two
global cocycle increments are preserved.  The averaged trace is
one-cell affine-periodic, and cellwise concavity gives
\begin{equation*}
m^{-2}\sum_ zK(c_ z)\leq K\!\left(m^{-2}\sum_ zc_ z\right).
\end{equation*}
Conversely such a trace repeats globally. This proves (18)--(19) on a
torus; the explicit finite-square affine-periodicity correction used later is
proved in Lemma~\ref{lem:two-scale-scalarization}.  Equations (15)--(16) say
exactly that $H=\mathcal H_ {h,L}\Psi$.  This identifies the deterministic
cell functional in (18)--(19).

\end{proof}

\section{A quantitative two-scale comparison}\label{sec:two-scale}

We now remove one fine band $\eta$ from the transport problem by
replacing $\Psi$ with $\mathcal H_{h,L}\Psi$.  The remaining source is
the independent coarse field $\zeta$.  The previous section treats a
constant base flux, so we first need room to smooth the coarse problem,
freeze it on blocks, and join the resulting cell competitors.
With the outer blocks of the fine band normalized to side one, the
scales satisfy
\[
1\ll L\ll B\ll A\ll S\leq R:
\]
Here $L$ is the cell side used by $\mathcal H_{h,L}$, and $B$ is the
side of a block on which the smoothed flux or slope is approximated by a
constant.  Smoothing on the larger scale $A$ makes this approximation
accurate, while $A\ll S$ keeps its effect on the coarse source small.
The ambient square has side $R$.

\begin{lemma}[Two-scale comparison]\label{lem:two-scale-scalarization}
Let $\eta$ be one retained $h$-level dyadic band, normalized so that
its outer blocks have side one. Let $\zeta$ be an independent sum of
at most $T$ coarser Haar levels whose smallest parent side, in this
normalization, is the dyadic integer $S$.  Work on $Q_ R$, where
$R$ is a common multiple of $L,B,A,S$, and all dyadic parents are
aligned with its boundary.  For brevity write $Q:=Q_R$ throughout this
lemma.  Let $P$ be the sum of the scaled linear
right inverses (12a), so that
$-\operatorname {div}P=\zeta$ and $P\cdot n=0$. It satisfies
\begin{equation}
\sup_ {x\in Q_ R}\mathbb E|P(x)|\leq C\sqrt T,\qquad
\mathbb E\fint_ {\mathbb T_ {2R}^2}
|P^{\rm ref}(x+z)-P^{\rm ref}(x)|\,dx\leq C|z|/S,
\tag{22}
\end{equation}
for $|z|\leq A\ll S$, where $P^{\rm ref}$ is the covariant
reflection defined in (22b0) below.  The same estimate, without
reflection, holds on every aligned interior translate.
For $1\ll L\ll B\ll A\ll S$, $0<\varepsilon<1/10$, and a
$D_ 4$-invariant $\Psi\in\mathscr P_ M$, put
\begin{equation}
\delta=\varepsilon+\frac1{\varepsilon L}
+\frac{L}{\varepsilon B}+\frac{B}{\varepsilon A}+\frac AS.
\tag{23}
\end{equation}
Assume also
\begin{equation}
\varepsilon L\geq C,\qquad L+B^2/A\leq c\varepsilon B.
\tag{23a}
\end{equation}
Both inequalities hold for the schedule used below.
On this aligned square,
\begin{equation}
\begin{split}
\frac1{|Q|}\bigl|\mathbb EJ_ \Psi^Q(\eta+\zeta)
-\mathbb EJ_ {\mathcal H_ {h,L}\Psi}^Q(\zeta)\bigr|
\leq {}&C\delta\left(
1+\frac{\mathbb EJ_ {\mathcal H_ {h,L}\Psi}^Q(\zeta)}{|Q|}
+M+\sqrt T+\sqrt h\right)\\
&+C\sqrt{\frac{\log(C/\varepsilon)}L}.                            
\end{split}
\tag{24}
\label{eq:two-scale-scalarization}
\end{equation}
The estimate is uniform for every iterated profile occurring below; in
particular, there is no uncontrolled error from an infinite-volume limit.

\end{lemma}

\begin{proof}
Write $H=\mathcal H_{h,L}\Psi$.  For the upper bound, we start with a
flux for $J_H^Q(\zeta)$ and add fine-source corrections with zero normal
trace on $B$-blocks.  The local estimate (26) will compare their cost
with $H$ at the frozen flux.  For the lower bound, we start with a dual
potential for $J_H^Q(\zeta)$ and insert fine cell optimizers while
preserving common values on the $B$-block boundaries.

Both constructions require smoothing near the outer boundary.  We first
prove the coarse translation estimate and a reflection argument that
preserves zero normal trace, including for measure-valued fluxes.  The
terms in $\delta$ then arise from contraction, $L$-cell corridors,
affine-periodicity correction, $B$-block seams, and coarse translation,
in that order.  The separate square-root term in (24) is the Gaussian
concentration error from Lemma~\ref{lem:finite-trace-net}.

\smallskip\noindent
\underline{\textit{Step 1.}} Control translations of the coarse flux.

We use the following explicit right inverse in (22).
On a Haar parent $[0,r]^2$, for a mean-zero density $f$, set
\begin{equation}
\begin{aligned}
m(y)&=r^{-1}\int_0^r f(s,y)\,ds,\\
P_1(x,y)&=-\int_0^x(f(s,y)-m(y))\,ds,\\
P_2(x,y)&=-\int_0^y m(t)\,dt.
\end{aligned}
\tag{22a0}
\end{equation}
Then $-\operatorname {div}P=f$ and $P\cdot n=0$ on all four
faces.  For a normalized Haar detail on a parent of side $r$, every
pointwise Gaussian coefficient in (22a0) is $O(1)$, while the expected
total variation of $DP$ per unit area is $O(r^{-1})$.
Consequently
\begin{equation}
\mathbb E\fint|P_ j(\cdot+z)-P_ j|
\leq C\min\{1,|z|/r_ j\}.
\tag{22a}
\end{equation}
The fields $P_ j$ are independent and centred, so their pointwise
second moments add, giving the first estimate in (22).  Since the
scales $r_ j\geq S$ grow geometrically, summing (22a) gives the second.
The explicit normal component vanishes on every aligned parent face;
covariant reflection therefore creates no jump there, and the same
BV translation estimate holds on the reflected torus.

\smallskip\noindent
\underline{\textit{Step 2.}} Smooth without changing the boundary constraint.

An effective minimizing flux may have a singular part.  We therefore
smooth in the space of vector measures and preserve the divergence and
normal-trace constraints before estimating the energy.  A vector
Radon measure $\mu\in\mathcal M(\overline Q_ R;\mathbb R^2)$ has
source $\rho$ and zero outer normal trace precisely when
\begin{equation}
\int_ {\overline Q_ R}\nabla\phi\cdot d\mu
=\langle\rho,\phi\rangle
\quad\text{for every }\phi\in C^1(\overline Q_ R).
\tag{22b0}
\end{equation}
In particular, if $w=P\mathcal L^2+v$ and $P\mathcal L^2$ have
the same source and zero normal trace, then $v$ is divergence-free
and the left side of (22b0) is zero.

Identify $\mathbb T_ {2R}^2$ with $[-R,R)^2$.  For
$\epsilon=(\epsilon_ 1,\epsilon_ 2)\in\{\pm1\}^2$, set
$S_ \epsilon x=(\epsilon_ 1x_ 1,\epsilon_ 2x_ 2)$ and
$O_ \epsilon=\operatorname {diag}(\epsilon_ 1,\epsilon_ 2)$.  Define
the covariant reflection, including any boundary-supported tangential
part, by
\begin{equation}
\int_ {\mathbb T_ {2R}^2}\Phi(y)\cdot dv^{\rm ref}(y)
=\sum_ {\epsilon\in\{\pm1\}^2}\int_ {\overline Q_ R}
\Phi(S_ \epsilon x)\cdot O_ \epsilon\,dv(x).
\tag{22b1}
\end{equation}
For every smooth periodic scalar $\chi$, the chain rule and (22b0)
give
\begin{equation}
\int_ {\mathbb T_ {2R}^2}\nabla\chi\cdot dv^{\rm ref}
=\sum_ \epsilon\int_ {\overline Q_ R}
\nabla(\chi\circ S_ \epsilon)\cdot dv=0.
\tag{22b2}
\end{equation}
Thus the reflected measure remains divergence-free across the faces,
periodic seams, and corners.  For a density $P$, define
$P^{\rm ref}(S_ \epsilon x)=O_ \epsilon P(x)$.  The same calculation
reflects its scalar source evenly.  Thus
$(P\mathcal L^2+v)^{\rm ref}=P^{\rm ref}\mathcal L^2+v^{\rm ref}$.

Let $\varphi_ A$ be a nonnegative smooth radial probability kernel on
the torus, supported at scale $A<R/4$, and put
\begin{equation}
v_ A=(\varphi_ A*v^{\rm ref})|_ {Q_ R},\qquad
P_ A=(\varphi_ A*P^{\rm ref})|_ {Q_ R}.
\tag{22b3}
\end{equation}
Convolution preserves (22b2), and covariance under the reflections
forces the classical normal component of $v_ A$ to vanish on every
side.  Hence $v_ A$ is an exact smooth divergence-free zero-normal
field on $Q_ R$, and
\begin{equation}
\int_ {Q_ R}|Dv_ A|\leq CA^{-1}|v|(\overline Q_ R).
\tag{22b4}
\end{equation}

To control the energy of this smoothing, let $H$ be any convex profile with
$|z|\leq H(z)\leq|z|+M_ H$.  Its recession function is $|z|$, and
for $\mu=\mu^a\mathcal L^2+\mu^s$ define
\begin{equation*}
\mathcal F_ H(\mu;U)=\int_ UH(\mu^a)\,dx+|\mu^s|(\overline U).
\end{equation*}
The perspective
\begin{equation*}
\widehat H(t,z)=
\begin{cases}tH(z/t),&t>0,\\ |z|,&t=0,
\end{cases}
\end{equation*}
is convex, lower semicontinuous, and positively one-homogeneous.  If
$\lambda$ dominates $\mathcal L^2$ and $\mu$, and
$(a,b)=d(\mathcal L^2,\mu)/d\lambda$, then
\begin{equation}
\mathcal F_ H(\mu)=\int\widehat H(a,b)\,d\lambda.
\tag{22b5}
\end{equation}
Applying convexity pointwise to the convolution of the pair
$(\mathcal L^2,\mu)$, followed by Fubini, proves the measure-level
Jensen inequality
\begin{equation}
\mathcal F_ H(\varphi_ A*\mu;\mathbb T_ {2R}^2)
\leq\mathcal F_ H(\mu;\mathbb T_ {2R}^2).
\tag{22b6}
\end{equation}
This argument includes singular measures and does not presuppose an
$L^1$ minimizer.

Assume now that $H$ is invariant under the signed coordinate
reflections, as is every profile in the iteration.  The reflected
measure has energy at most four times the original energy (overlapping
face masses use the triangle inequality for the recession norm), while
the convolved density has four equal quadrant energies.  From (22b6),
\begin{equation}
\int_ {Q_ R}H(P_ A+v_ A)\,dx
\leq\mathcal F_ H(P\mathcal L^2+v;Q_ R).
\tag{22b7}
\end{equation}
Moreover (22) and Fubini give
\begin{equation}
\mathbb E\int_ {Q_ R}|P_ A-P|\leq C(A/S)|Q_ R|.
\tag{22b8}
\end{equation}
Since $H$ is $1$-Lipschitz, (22b7)--(22b8) yield
\begin{equation}
\mathbb E\int_ {Q_ R}H(P+v_ A)
\leq\mathbb E\mathcal F_ H(P\mathcal L^2+v;Q_ R)
+C(A/S)|Q_ R|.
\tag{22b9}
\end{equation}
Consequently, we may smooth the divergence-free correction on the
whole square at expected cost $CA/S$ per unit area.

\smallskip\noindent
\underline{\textit{Step 3.}} Upper bound: smooth and freeze the flux.

For the upper bound in \eqref{eq:two-scale-scalarization}, use the relaxed linear-growth functional on
vector Radon measures and take a near-minimizer
$w=P\mathcal L^2+v$ for
$J_ {\mathcal H_ {h,L}\Psi}^Q(\zeta)$, where
$\operatorname {div}v=0$ and $v$ has zero outer normal trace.
(The singular part is charged by the recession
function $|\cdot|$.) Convolve only $v$ on scale $A$.  Equations
(22b7)--(22b9) change the effective energy by at most $CA/S$ per unit
area.  Coercivity
$\mathcal H_ {h,L}\Psi\geq|\cdot|$ controls the total variation of $v$.
Let $p_ D=\fint_ D(P+v_ A)$ on each interior $B$-block.  Poincare's
inequality, (22), and (22b4) give
\begin{equation}
R^{-2}\mathbb E\sum_ D\int_ D|P+v_ A-p_ D|
\leq C\left(\frac BA+\frac BS\right)
\left(1+R^{-2}\mathbb E|v|(Q_ R)+\sqrt T\right).
\tag{25}
\end{equation}

The remaining local task is to insert the fine source at constant base
$p_D$.  We next prove a deterministic comparison uniform in this base,
so that it can subsequently be applied conditionally to the coarse field.

\smallskip\noindent
\underline{\textit{Step 4.}} Upper bound: compare the finite-block cost.

We seek the two inequalities in (26) between the expected cost of a
zero-normal fine correction on a $B$-box and $H(p)$.  Repeating an
optimal affine-periodic $L$-cell trace gives the lower inequality.  For
the upper inequality, the finite-trace lemma first makes the trace
deterministic.  Averaging that trace over cell translations leaves small
errors in its opposite-face increments; correcting those errors will
produce an admissible trace for $H$.

The fine source is exactly balanced on every $B$-block.  To formulate
the local problem without suppressing singular parts, for $p\in
\mathbb R^2$ set
\begin{equation}
\begin{split}
I_ B(p,\eta_ D)&=\min\{\mathcal F_ \Psi(p\mathcal L^2+u;D):
-\operatorname {div}u=\eta_ D, [u\cdot n]_ {\partial D}=0\},\\
N_ B(p)&=B^{-2}\mathbb E_ \eta I_ B(p,\eta_ D).                    
\end{split}
\tag{25a0}
\end{equation}
Here the constraint is on the correction $u$, not on the constant
base $p\mathcal L^2$.  Apply the same separation argument as in (2)
to this affine constraint.  Fenchel's inequality reads
$\Psi(p+u)\geq\nabla F\cdot(p+u)-\Psi^*(\nabla F)$, while the
zero-normal constraint gives
$\int_ D\nabla F\cdot du=\langle\eta_ D,F\rangle$.  Separation gives
the reverse inequality and hence, realization by realization,
\begin{equation}
I_ B(p,\eta_ D)=\max_ {F\in\mathscr A_ B}
\left\{\langle\eta_ D,F\rangle+
\int_ Dp\cdot\nabla F-\int_ D\Psi^*(\nabla F)\right\}.
\tag{25a1}
\end{equation}

We first prove the lower side of the finite comparison.  Fix
$|a|\leq1$ and a trace $c\in\mathcal T_ L^{\rm per}(a)$ attaining
$\mathcal R(a)$.  Repeat its universal-cover affine-periodic lift through the $B$-box
and, for each fine realization, fill every $L$-cell by an optimizer
with that complete trace.  The interiors paste, and every cell has mean
gradient $a$.  Substitution in (25a1), expectation, and maximization
over $a$ give
\begin{equation}
(\mathcal H_ {h,L}\Psi)(p)\leq N_ B(p).
\tag{25a2}
\end{equation}

For the reverse side apply (21a')--(21e) to (25a1) with the constant
realized field $P\equiv p$.  Since (21d) applies exactly,
\begin{align}
N_ B(p)&\leq \mathfrak A_ {B,L}(p)
+C(M+\sqrt h+|p|)
\left(\varepsilon+\frac1{\varepsilon L}\right)
+C\sqrt{\frac{\log(C/\varepsilon)}L},                         \tag{25a3}\\
\mathfrak A_ {B,L}(p)&=B^{-2}\max_ {c\in\mathfrak T_ {B,L}}
\sum_ z\{K(c_ z)+p\cdot m_ L(c_ z)\}.                       \tag{25a4}
\end{align}
Every constant is uniform in the possibly unbounded vector $p$.

Choose a maximizing trace in (25a4), one global extension $F$, and
write $n=B/L$,
\begin{equation}
g(s)=n^{-2}\sum_ {z\in\{0,\ldots,n-1\}^2}F(Lz+s),\quad s\in Q_ L,
\qquad q=B^{-2}\int_ {Q_ B}\nabla F.
\tag{25a5}
\end{equation}
Additive constants do not affect $K$, so concavity and the divergence
theorem give
\begin{equation}
n^{-2}\sum_ zK(c_ z)\leq K(g|_ {\partial Q_ L}),\qquad
B^{-2}\sum_ zp\cdot m_ L(c_ z)=p\cdot q.
\tag{25a6}
\end{equation}
For $s$ on the face $s_ j=0$, let
$D_ j(s)=g(s+Le_ j)-g(s)$.  Telescoping in direction $e_ j$ leaves
only the two outer rows of the $B$-box.  The resulting function is
$2/n$-Lipschitz and its face average is $Lq_ j$; hence
\begin{equation}
\|D_ j-Lq_ j\|_ \infty\leq2L/n.
\tag{25a7}
\end{equation}

The averaged function $g$ is already Lipschitz, but (25a7) gives only
approximate affine-periodicity.  We repeat $g$ with the desired affine
increments and then remove the resulting small seam jumps by a Lipschitz
envelope.  Put
$\theta=1-4\varepsilon$, $d=2L/n$, and, on half-open cells of the
universal cover, define
\begin{equation*}
G_ 0(s+Lk)=\theta g(s)+\theta Lq\cdot k.
\end{equation*}
Its seam jumps are at most $d$.  A segment of Euclidean length $r$
crosses at most $2+\sqrt2r/L$ grid faces: the sum of its two coordinate
displacements is at most $\sqrt2r$.  Since
$\sqrt2d/L=2\sqrt2/n\leq\varepsilon$ when $n\varepsilon\geq8$, we obtain
\begin{equation}
|G_ 0(x)-G_ 0(y)|\leq(1-3\varepsilon)|x-y|+4d.
\tag{25a8}
\end{equation}
Since $|\theta q|\leq\theta<1-\varepsilon/2$, the following
infimum is finite.  Define
\begin{equation*}
\widehat G(x)=\inf_ y\{G_ 0(y)+(1-\varepsilon/2)|x-y|\}.
\end{equation*}
The change of variable $y\mapsto y+Lk$ proves the exact cocycle
$\widehat G(x+Lk)=\widehat G(x)+\theta Lq\cdot k$, while (25a8)
gives $G_ 0-4d\leq\widehat G\leq G_ 0$.  Thus
$\widehat c=\widehat G|_ {\partial Q_ L}$ is an admissible
affine-periodic trace of slope $\theta q$ and differs from
$\theta g$ by at most $4d$.

For a fine realization take an optimizer with trace $g$, keep its
$\theta$-multiple on every unit outer block farther than
$C(1+d/\varepsilon)$ from $\partial Q_ L$, prescribe
$\widehat c$ on the boundary, and use McShane extension in the collar.
The same realized flux calculation as (21a') and
$-\Psi^*(\theta a)\geq-\theta\Psi^*(a)
-(1-\theta)\Psi^*(0)$ give
\begin{equation}
L^{-2}K(g)\leq \mathcal R(\theta q)+C(M+\sqrt h)
\left(\varepsilon+\frac1{\varepsilon L}
+\frac L{\varepsilon B}\right).
\tag{25b}
\end{equation}
Since $p\cdot q\leq p\cdot(\theta q)+4\varepsilon|p|$, equations
(25a3)--(25b) and the definition of $\mathcal H$ prove the complete
finite-block zero-normal comparison
\begin{equation}
\boxed{\begin{aligned}
(\mathcal H_ {h,L}\Psi)(p)\leq N_ B(p)
\leq{}&(\mathcal H_ {h,L}\Psi)(p)\\
&+C(M+\sqrt h+|p|)
\left(\varepsilon+\frac1{\varepsilon L}
+\frac L{\varepsilon B}\right)
+C\sqrt{\frac{\log(C/\varepsilon)}L}.
\end{aligned}}
\tag{26}
\end{equation}

Estimate (26) is uniform in the possibly unbounded base vector $p$,
so it applies conditionally to the random block averages $p_D$.

\smallskip\noindent
\underline{\textit{Step 5.}} Upper bound: paste the local flux corrections.

Now condition on the coarse sigma-field.  Each
$p_ D=\fint_ D(P+v_ A)$ is then deterministic and independent of the
fine field, and the local sources are independent.  Using the same
correction $u$ at two bases changes its normalized energy by at most
$|p-p'|$, so $N_ B$ is $1$-Lipschitz and hence Borel.  Equations
(25a0) and (26) give the exact
conditional calculation
\begin{equation}
\begin{split}
\mathbb E_ \eta\!\left[\sum_ DI_ B(p_ D,\eta_ D)\mid\mathscr G\right]
={}&\sum_ D|D|N_ B(p_ D)\\
\leq{}&\sum_ D|D|H(p_ D)
+C\alpha\left\{(M+\sqrt h)|Q|+\sum_ D|D||p_ D|\right\}\\
&+C|Q|\sqrt{\frac{\log(C/\varepsilon)}L},\qquad
\alpha=\varepsilon+\frac1{\varepsilon L}
+\frac L{\varepsilon B},               
\end{split}
\tag{26i}
\end{equation}
where $H=\mathcal H_ {h,L}\Psi$.  We now use this conditional bound to construct a global flux and then
average over the coarse field.

Fix $\tau_{\rm sel}>0$ and choose measurable
$\tau_{\rm sel}|D|$-minimizers $u_ D^{\tau_{\rm sel}}$ in (25a0)
as justified at the end of the proof.  The zero extensions have no interface divergence.  Add them to the
unfrozen smoothed flux to obtain the global competitor
\begin{equation*}
\mu^{\tau_{\rm sel}}=(P+v_ A)\mathcal L^2+\sum_ Du_ D^{\tau_{\rm sel}},
\end{equation*}
which has source $\eta+\zeta$ and zero outer normal trace.  The constants
$p_D$ enter only the energy comparison below.  The $1$-Lipschitz property of $\Psi$ and
subadditivity of its recession norm give
\begin{equation}
\mathcal F_ \Psi(\mu^{\tau_{\rm sel}};Q)
\leq\sum_ D\mathcal F_ \Psi(p_ D\mathcal L^2+u_ D^{\tau_{\rm sel}};D)
+\sum_ D\int_ D|P+v_ A-p_ D|.
\tag{26j}
\end{equation}
This remains true if local corrections have tangential mass on a common
face.  Take conditional expectation, use (26i), Jensen
$\sum_ D|D|H(p_ D)\leq\int_ QH(P+v_ A)$, and
$\sum_ D|D||p_ D|\leq\int_ QH(P+v_ A)$.  Equations (22b9) and (25),
followed by $\tau_{\rm sel}\downarrow0$, prove the upper inequality in
\eqref{eq:two-scale-scalarization}.
Thus
$\mathbb EJ_ \Psi(\eta+\zeta)\leq\mathbb EJ_ H(\zeta)+\text{error}$,
with no additional interface cost from the local corrections.

\smallskip\noindent
\underline{\textit{Step 6.}} Lower bound: smooth the potential and freeze its slopes.

For the lower bound in \eqref{eq:two-scale-scalarization}, fix
$\tau_{\rm sel}>0$, use (2) for
$J_ {\mathcal H_ {h,L}\Psi}^Q(\zeta)$, and take a
$\tau_{\rm sel}|Q|$-maximizing global potential $F$, normalized at one point.
Extend it by successive even
reflections across the square sides and then $2R$-periodically, and
convolve at scale $A$.  Even reflection and convolution preserve the
unit Lipschitz bound: explicitly, the coordinate-folding map from the
torus to $Q_ R$ is $1$-Lipschitz.  If
$F_ A=\varphi_ A*F^{\rm ref}$, then
$\|D^2F_ A\|_ \infty\leq C/A$.  Reflecting $P$ covariantly as in
(22b1), Fubini gives the exact identity
\begin{equation*}
\int_ {\mathbb T_ {2R}^2}P^{\rm ref}\cdot\nabla F_ A
=\int_ {\mathbb T_ {2R}^2}P_ A\cdot\nabla F^{\rm ref}.
\end{equation*}
For $\mathcal R=-H^*$, concavity and reflection invariance give
$\mathcal R(\nabla F_ A)\geq\varphi_ A*\mathcal R(\nabla F^{\rm ref})$.  Divide the
resulting torus inequality by four and use (22b8).  The value decreases
by at most $CA|Q|/S$, and no outer collar is discarded.

Write $H=\mathcal H_ {h,L}\Psi$, $\mathcal R=-H^*$, and
$M_ h=M+C\sqrt h$.  Concavity of $\mathcal R$, even reflection, and (22) give
for the smoothed potential
\begin{equation}
\mathbb E\int_ Q\{P\cdot\nabla F_ A+\mathcal R(\nabla F_ A)\}
\geq\mathbb EJ_ H^Q(\zeta)-C(A/S+\tau_{\rm sel})|Q|,
\tag{26a}
\end{equation}
where $\tau_{\rm sel}|Q|$ is the initial near-maximization error.  Also
$\|D^2F_ A\|_ \infty\leq C/A$.

For every $B$-block $D$, put
\begin{equation*}
a_ D=\fint_ D\nabla F_ A.
\end{equation*}
Then $|a_ D|\leq1$, and Jensen's inequality for the concave function
$\mathcal R$ gives
\begin{equation}
|D|\mathcal R(a_ D)\geq\int_ D \mathcal R(\nabla F_ A).
\tag{26b}
\end{equation}
Choose measurably a trace $c_ D\in\mathcal T_ L^{\rm per}(a_ D)$
attaining $\mathcal R(a_ D)$, and repeat one of its cell extensions as
$a_ D\cdot x+u_ D(x)$.  The periodic part may be normalized so that
$\|u_ D\|_ \infty\leq CL$.  Let $\ell_ D$ be affine with gradient
$a_ D$ and the value of $F_ A$ at the centre of $D$, and put
$f_ D=\ell_ D+u_ D$, with phase aligned to the global $L$-grid.
Since $\|D^2F_ A\|_ \infty\leq C/A$,
\begin{equation}
\|F_ A-f_ D\|_ {L^\infty(D)}\leq C(L+B^2/A).
\tag{26c}
\end{equation}

Estimate (26c) controls the error in values, which is what we need to
join the block potentials while keeping the Lipschitz constraint.

\smallskip\noindent
\underline{\textit{Step 7.}} Lower bound: join the fine potentials.

The potentials $f_D$ may disagree on adjacent block boundaries.
We prescribe the common value $\theta F_A$ there and interpolate to
$\theta f_D$ inside each block.  The contraction provides the slack
needed for the interpolation gradient; away from these seams, the
affine-periodic traces remain available for fine optimization.

Set $\theta=1-4\varepsilon$ and
\begin{equation*}
w=C(L+B^2/A)/\varepsilon<B/4
\end{equation*}
by (23a).  Choose a cutoff $\chi_ D$ which is zero on $\partial D$,
one outside the width-$w$ seam, and satisfies
$|\nabla\chi_ D|\leq C/w$.  Define
\begin{equation}
G_ D=\theta\{(1-\chi_ D)F_ A+\chi_ D f_ D\}.
\tag{26d}
\end{equation}
The first part of its gradient is $\theta$ times a convex combination
of two vectors in $\overline B_ 1$, while (26c) makes the cutoff term
at most $\varepsilon$.  Hence $|\nabla G_ D|\leq1-3\varepsilon$.
Because $G_ D=\theta F_ A$ on every $\partial D$, the $G_ D$'s paste
into one global Lipschitz skeleton.  The proportion of $L$-cells
meeting a transition seam is at most
\begin{equation}
C\left(\frac{L}{\varepsilon B}
+\frac{B}{\varepsilon A}\right).
\tag{26e}
\end{equation}

In a non-seam $L$-cell the trace is, up to an additive constant, a
translate of $\theta c_ D$.  Concavity of $K$, together with
$K(0)\geq0$ from (16d), gives
\begin{equation}
K(\theta c_ D)\geq\theta K(c_ D)=\theta L^2\mathcal R(a_ D).
\tag{26f}
\end{equation}
Let $t_ C=G_ D|_ {\partial C}$ be the trace on each $L$-cell and,
conditionally on the coarse sigma-field, fill it by a measurable
realization-wise optimizer for the independent fine Gaussian vector in
that cell.  Complete traces make the interiors paste exactly to a global
potential $\widetilde F$.  The fixed-trace identity (21c) now gives,
without interchanging a random maximum and expectation,
\begin{equation}
\mathbb E_ \eta\!\left[
\langle\eta,\widetilde F\rangle-
\int_ Q\Psi^*(\nabla\widetilde F)\,\middle|\,\mathscr G\right]
=\sum_ CK(t_ C).
\tag{26g}
\end{equation}
In a seam cell $K(t_ C)\geq0$: use the deterministic trace extension,
then average the centred source, and use $-\Psi^*\geq0$.  Therefore
(26b), (26e), and (26f), with $0\leq \mathcal R\leq M_ h$, imply
\begin{equation}
\sum_ CK(t_ C)\geq\theta\int_ Q\mathcal R(\nabla F_ A)
-CM_ h|Q|\left(\varepsilon+\frac{L}{\varepsilon B}
+\frac{B}{\varepsilon A}\right).
\tag{26h}
\end{equation}

It remains to control the coarse-source pairing after replacing the
interiors.  Put $P_ D=\fint_ DP$.  The potentials
$\widetilde F$ and $\theta F_ A$ have identical traces on
$\partial D$, and hence
\begin{equation}
\int_ DP_ D\cdot\nabla(\widetilde F-\theta F_ A)=0,
\qquad
\int_ DP\cdot\nabla\widetilde F
\geq\theta\int_ DP\cdot\nabla F_ A-2\int_ D|P-P_ D|.
\tag{26k}
\end{equation}
Averaging (22) over pairs in each $B$-block gives
\begin{equation}
\mathbb E\sum_ D\int_ D|P-P_ D|\leq C(B/S)|Q|.
\tag{26l}
\end{equation}
Finally, replacing the effective value by its factor $\theta$ costs
at most $C\varepsilon(M_ h+\sqrt T)|Q|$, because
$0\leq \mathcal R\leq M_ h$ and $\sup_ x\mathbb E|P(x)|\leq C\sqrt T$.
Combine (26a), (26g)--(26l), and let $\tau_{\rm sel}\downarrow0$.  This proves
the lower inequality in \eqref{eq:two-scale-scalarization} with exactly the terms in (23), and establishes
the direction
$\mathbb EJ_ \Psi(\eta+\zeta)\geq\mathbb EJ_ H(\zeta)-\text{error}$.

The constructions above use choices depending on the coarse field or
on individual fine-source realizations.  We finish by verifying the
measurability needed for their conditional expectations.

\smallskip\noindent
\underline{\textit{Step 8.}} Measurable selections.

We use the following selection principle.  If $S$ is standard Borel, $X$ is a
countable union of compact metrizable sets, $\Gamma\subset S\times X$
is Borel with nonempty closed sections in the relevant compact sets, and
$f$ is jointly Borel, lower semicontinuous in $x$, and has compact
parameterwise sublevels, then the infimum is universally measurable and
there is a universally measurable $\tau_{\rm sel}$-minimizer for every
$\tau_{\rm sel}>0$.  Apply this to $-f$ for upper semicontinuous maxima.
Indeed, intersect rational sublevel sets with the compact exhaustion and
select successively the first member of a countable finite-ball cover
whose section is nonempty; the nested balls give the selector.  On our
completed Gaussian spaces universal measurability is sufficient.

For a closed block $D$, take
$X_ D=\mathcal M(\overline D;\mathbb R^2)$ with its weak-* Borel
structure.  Since $C(\overline D;\mathbb R^2)$ is separable, every
total-variation ball is compact metrizable, and these balls exhaust
$X_ D$.  Choose a countable $C^1$-dense family $\{\phi_ k\}$.  If
the finite Gaussian vector $g$ parametrizes $\eta_ D(g)$, the graph
\begin{equation}
\Gamma_ D=\left\{(g,u):
\int_ {\overline D}\nabla\phi_ k\cdot du
=\langle\eta_ D(g),\phi_ k\rangle\quad\text{for every }k\right\}
\tag{26m}
\end{equation}
is closed and is exactly the divergence plus zero-normal constraint by
density.  The function
\begin{equation}
(p,g,u)\longmapsto\mathcal F_ \Psi(p\mathcal L^2+u;D)
\tag{26n}
\end{equation}
is jointly Borel and weak-* lower semicontinuous by the perspective
formula (22b5).  Its sublevels are compact because
\begin{equation}
|u|(\overline D)\leq
\mathcal F_ \Psi(p\mathcal L^2+u;D)+|p||D|.
\tag{26o}
\end{equation}
The explicit zero-normal Haar right inverse makes every constraint
section nonempty.  The principle therefore supplies a jointly
universally measurable $u_ D^{\tau_{\rm sel}}(p,g)$ satisfying the constraints
and
\begin{equation}
\mathcal F_ \Psi(p\mathcal L^2+u_ D^{\tau_{\rm sel}};D)
\leq I_ B(p,\eta_ D(g))+\tau_{\rm sel}|D|.
\tag{26p}
\end{equation}
This proves the measurability asserted in (26j), including composition
with an unbounded coarse-measurable $p_ D$.

For the effective primal problem use the same measure space on $Q$,
the closed subspace
\begin{equation*}
\mathscr V_ Q=\left\{v:\int_ {\overline Q}\nabla\phi_ k\cdot dv=0
\text{ for every }k\right\},
\end{equation*}
and parameter space $L^1(Q;\mathbb R^2)$ for $P$.  The map
$(P,v)\mapsto\mathcal F_ H(P\mathcal L^2+v;Q)$ is jointly Borel,
weak-* lower semicontinuous in $v$, and coercive exactly as in
(26n)--(26o).  We obtain a universally measurable
$\tau_{\rm sel}|Q|$-minimizer $v^{\tau_{\rm sel}}(P)$.  Reflection, convolution,
block averaging, and finite sums are Borel operations, so the entire
competitor for the upper bound is measurable.

For dual selections, the spaces
\begin{equation*}
\mathscr A_ L=\{F\in C(\overline Q_ L):F(0)=0,
\operatorname {Lip}(F)\leq1\},
\qquad \mathcal T_ L=\{F|_ {\partial Q_ L}:F\in\mathscr A_ L\}
\end{equation*}
are compact Polish in the uniform norm, and the fixed-trace graph
$\{(c,F):F|_ {\partial Q_ L}=c\}$ is closed.  Uniform convergence in
$\mathscr A_ L$ implies weak-* convergence of gradients.  Hence, for
the cell Gaussian vector $g$,
\begin{equation}
(g,F)\longmapsto\langle\eta(g),F\rangle
-\int_ {Q_ L}\Psi^*(\nabla F)
\tag{26q}
\end{equation}
is jointly Borel and upper semicontinuous in $F$.  The maximum
principle gives a jointly measurable exact optimizer
$F^{\rm opt}(c,g)$.  For an unanchored trace, subtract its value at a
fixed corner, apply the selector, and add the value back.  This proves
the exact conditional identity (26g), not merely an inequality between
outer expectations.

The full affine-periodic trace graph
\begin{equation}
\mathfrak G_ L=\{(a,c)\in\overline B_ 1\times\mathcal T_ L:
c(L,t)-c(0,t)=La_ 1,\ c(t,L)-c(t,0)=La_ 2\}
\tag{26r}
\end{equation}
is compact.  The kernel $K(c)$ is upper semicontinuous: use the
upper semicontinuity of the realization-wise maximum and the pathwise,
trace-independent domination
\begin{equation*}
K^{\rm raw}(c,\eta)\leq\int_ {Q_ L}(|Q|+M),
\end{equation*}
where
$K^{\rm raw}$ denotes the maximum before expectation, in reverse
Fatou.  Its maximizing sections are
nonempty compact, so there is a measurable
$a\mapsto c^*(a)$ with $K(c^*(a))=L^2\mathcal R(a)$.  A measurable periodic
interior extension is supplied explicitly by
\begin{equation*}
E(c)(x)=\inf_ {y\in\partial Q_ L}\{c(y)+|x-y|\}
\end{equation*}
and affine-cocycle repetition; $c\mapsto E(c)$ is Lipschitz in the
uniform norm on compact sets.  This justifies $c_ D,f_ D$, and every
cutoff in the lower bound.  For the strict-slope extension in Definition~\ref{def:cell-operator}, one does not use the McShane formula.  Instead restrict
the potential correspondence to $\nabla F\in K_ \Psi$; weak-* closure
of the closed convex gradient constraint gives a compact closed graph,
and the same selection principle supplies a compatible interior.

Finally, on the compact Polish space $\mathscr A_ Q$, the effective
dual objective
\begin{equation}
(P,F)\longmapsto\int_ QP\cdot\nabla F+\int_ Q\mathcal R(\nabla F)
\tag{26s}
\end{equation}
is jointly Borel and upper semicontinuous in $F$: the first term is
continuous for $P\in L^1$, and $-\mathcal R$ is convex lower
semicontinuous.  It therefore has a coarse-measurable maximizing (or
$\tau_{\rm sel}$-maximizing) potential.  Thus all conditional expectations in
(21c), (26i), and (26g) are legitimate.  Letting the deterministic
$\tau_{\rm sel}\downarrow0$ in the two
competitor constructions proves \eqref{eq:two-scale-scalarization}.

\end{proof}

\section{Analytic consistency estimates}\label{sec:consistency}

To identify the limit of the iterated profiles, we must determine the
change caused by one group of $h$ levels.  For a profile viewed on the
spatial scale $\sqrt T$, a local Hessian $H/\sqrt T$ produces a
quadratic transport problem with energy of order $h/\sqrt T$.
Lemma~\ref{lem:grouped-quadratic-generator} computes its coefficient
and supplies the moment bounds needed to control the nonlinear errors.

The two inequalities require different uses of this quadratic problem.
For the upper bound, a Neumann flux is an admissible primal competitor;
its Gaussian tails allow us to use only a local expansion of the
profile.  For the lower bound, a Dirichlet potential must remain inside
the region where the dual cost is quadratic.  We first prove the
zero-boundary truncation needed to make that potential admissible.
Finally, Lemma~\ref{lem:global-barriers} turns a local lower contact
into an order relation on the whole space, so that the lower estimate
can be used with monotonicity of the cell operator.

\subsection{Lipschitz truncation with zero boundary data}

The truncation below is used in Lemma~\ref{lem:nonlinear-consistency}.
Besides controlling the gradient, it preserves the zero boundary values
that allow Dirichlet potentials on different blocks to be pasted.

\begin{lemma}\label{lem:lipschitz-truncation}
Let $D$ be a square or flat torus and let $g\in W^{1,4}(D)$; in the
square case assume $g\in W^{1,4}_ 0(D)$.  For every $\lambda>0$ there
are an open set $E\subset D$ and a function $g_ \lambda$ such that, in
the square case,
\begin{equation*}
g_ \lambda\in W^{1,\infty}(D)\cap W^{1,4}_ 0(D)
\quad\hbox{and its zero extension is globally Lipschitz},
\end{equation*}
and in either case
\begin{equation}
\|\nabla g_ \lambda\|_ \infty\leq C\lambda,\qquad
g_ \lambda=g\quad\hbox{on }D\setminus E,
\tag{27}
\end{equation}
and
\begin{equation}
|E|\leq C\lambda^{-4}\int_ D|\nabla g|^4,
\qquad
\int_ D|\nabla g_ \lambda-\nabla g|^2
\leq C\lambda^{-2}\int_ D|\nabla g|^4.
\tag{28}
\end{equation}

\end{lemma}

\begin{proof}
We keep $g$ where its gradient maximal function is at most $\lambda$
and extend these values with a controlled Lipschitz constant.  For a
square, the extension must also take the prescribed value zero at every
exterior point.  We verify this compatibility before constructing the
extension; the maximal estimate will then control both errors in (28).

\smallskip\noindent
\underline{\textit{Step 1.}} Lipschitz control on the good set.

In the boundary case let $\widetilde g$ be the zero extension of $g$
to $\mathbb R^2$.  Since $g\in W^{1,4}_ 0(D)$,
$\widetilde g\in W^{1,4}(\mathbb R^2)$ and
$\nabla\widetilde g=1_ D\nabla g$.  Set
\begin{equation*}
\mathcal O=\{x\in\mathbb R^2:
\mathcal M(|\nabla\widetilde g|)(x)>\lambda\},\qquad
E=D\cap\mathcal O.
\end{equation*}
At Lebesgue points of any Sobolev function $v$, telescoping its averages
over the balls $B(x,2^{-k}|x-y|)$, and then comparing the averages at
$x$ and $y$ on one ball of radius $3|x-y|$, gives
\begin{equation*}
|v(x)-v(y)|\leq C|x-y|
\{\mathcal M(|\nabla v|)(x)+\mathcal M(|\nabla v|)(y)\}.
\end{equation*}
Consequently, if $x,y\in D\setminus E$,
\begin{equation}
|g(x)-g(y)|\leq C\lambda|x-y|.
\tag{28a}
\end{equation}
To include zero exterior data in the extension, we also need a
Lipschitz bound between a good interior point and a point outside $D$.

If $x\in D$ and $d=\operatorname {dist}(x,D^c)$, then
$B(x,2d)\setminus D$ occupies a fixed fraction of $B(x,2d)$.
Telescoping from $x$ to the average on that ball, followed by
$L^1$-Poincare and $\widetilde g=0$ on the exterior fraction, gives
\begin{equation*}
|\widetilde g(x)|\leq Cd\mathcal M(|\nabla\widetilde g|)(x).
\end{equation*}
Thus, if also $x\notin\mathcal O$ and $y\in D^c$,
\begin{equation}
|g(x)|\leq C\lambda d\leq C\lambda|x-y|.
\tag{28b}
\end{equation}

\smallskip\noindent
\underline{\textit{Step 2.}} Extension with exact zero trace.

It follows that the prescribed data
\begin{equation*}
g\ \hbox{on }D\setminus\mathcal O,\qquad
0\ \hbox{on }\mathbb R^2\setminus D
\end{equation*}
are $C\lambda$-Lipschitz on their union.  Let
$A=(D\setminus\mathcal O)\cup D^c$, let $f$ denote these data, and
take the McShane extension
\begin{equation*}
\widetilde g_ \lambda(x)=\inf_ {a\in A}
\{f(a)+C\lambda|x-a|\}.
\end{equation*}
Every exterior point was prescribed zero, so
$\widetilde g_ \lambda=0$ on $\mathbb R^2\setminus D$.
Consequently $g_ \lambda=\widetilde g_ \lambda|_ D$ belongs to
$W^{1,\infty}(D)$, agrees with $g$ off $E$, and satisfies
$\|\nabla g_ \lambda\|_ \infty\leq C\lambda$.  Moreover
$|g_ \lambda(x)|\leq C\lambda\operatorname {dist}(x,D^c)$.
If $\chi_ k$ is zero at distance $<1/k$, one at distance $>2/k$,
and $|\nabla\chi_ k|\leq Ck$, then
$\chi_ kg_ \lambda\to g_ \lambda$ in $W^{1,4}(D)$: both gradient-error
terms are bounded and supported on a boundary strip whose measure tends
to zero, while $|g_ \lambda\nabla\chi_ k|\leq C\lambda$ there.  Hence
after mollifying each compactly supported Lipschitz cutoff,
$g_ \lambda\in W^{1,4}_ 0(D)$.  The required zero-trace property is
therefore in $W^{1,4}$; no closure assertion in the
$W^{1,\infty}$-norm is needed.

\smallskip\noindent
\underline{\textit{Step 3.}} Exceptional-set and gradient estimates.

The strong $L^4$ maximal estimate followed by Chebyshev gives the
desired measure bound:
\begin{equation}
|E|\leq C\lambda^{-4}\int_ D|\nabla g|^4.
\tag{28c}
\end{equation}
Since the gradients agree off $E$,
\begin{equation*}
\begin{split}
\int_ D|\nabla g_ \lambda-\nabla g|^2
&\leq C\lambda^2|E|+C\int_ E|\nabla g|^2\\
&\leq C\lambda^{-2}\int_ D|\nabla g|^4,
\end{split}
\end{equation*}
where the last step is Cauchy--Schwarz followed by (28c).  On a flat
torus use the periodic maximal function and the periodic McShane
extension of the good-set data; the same calculation proves (27)--(28).

\end{proof}

\subsection{Quadratic energy and Gaussian moments}

The quadratic cost with matrix $H$ has dual cost with matrix $H^{-1}$.
Its optimizer is therefore obtained by solving the elliptic equation
with operator $L_H$ below.  In dimension two, the inverse eigenvalues
accumulate logarithmically up to the cell resolution $2^{-h}$; their
coefficient will be the generator of the limiting equation.  We also
need this coefficient to be independent of the boundary trace and of
the period $L$, and need moments of the optimizers to estimate the
nonlinear remainders.

\begin{lemma}\label{lem:grouped-quadratic-generator}
Let $H$ be a positive definite $2\times2$ matrix and
$n=2^h$.  On the unit square let $\Pi_ h$ be orthogonal projection
onto functions constant on the $n^2$ small squares, followed by
removal of the constants, and put
\begin{equation*}
L_ H=-\operatorname {div}(H^{-1}\nabla).
\end{equation*}
If $E_ {h,H}^N$ is the expected minimum of
$\frac12\int u^THu$ under
$-\operatorname {div}u=\eta_ h$ and zero normal boundary condition,
and $E_ {h,H}^D$ is the corresponding Dirichlet dual value, then
\begin{equation}
E_ {h,H}^N=\frac12\operatorname {Tr}(\Pi_ h(L_ H^N)^{-1}),
\qquad
E_ {h,H}^D=\frac12\operatorname {Tr}(\Pi_ h(L_ H^D)^{-1}),
\tag{29}
\end{equation}
and
\begin{equation}
E_ {h,H}^D,E_ {h,H}^N
=\frac{h\log2}{4\pi}\sqrt{\det H}+o(h).
\tag{30}
\end{equation}
The error is uniform for $H$ in compact subsets of the positive
definite cone.  More precisely, define the deterministic full-trace
quadratic value
\begin{equation}
\mathcal E^{\rm tr}_ {h,L}(H)=L^{-2}\max_ c\mathbb E
\sup_ {f|_ {\partial Q_ L}=c}
\left\{\int f\eta_ h^{(L)}
-\frac12\int\nabla f^TH^{-1}\nabla f\right\},
\tag{30a}
\end{equation}
where $c$ ranges over all affine-periodic complete traces (the affine
slope is optimized as part of $c$).  Then
\begin{equation}
E_ {h,H}^D\leq\mathcal E^{\rm tr}_ {h,L}(H)\leq E_ {h,H}^N,
\tag{30b}
\end{equation}
so it has the leading term (30), uniformly in the auxiliary period
$L$.

The Gaussian primal and dual optimizers used in (29) also satisfy
\begin{equation}
\mathbb E\fint |u|^r\leq C_ {H,r}h^{r/2},\qquad r=2,3,4.
\tag{31}
\end{equation}

\end{lemma}

\begin{proof}
We first express each expected quadratic energy as a trace and compute
its logarithmic leading term.  The Dirichlet and Neumann values then
bound the complete-trace value, uniformly in $L$.  The remaining task
is (31): a pointwise variance estimate for the Gaussian optimizers
will imply all three moment bounds.

\smallskip\noindent
\underline{\textit{Step 1.}} Variational trace identities.

For the Neumann problem, introduce a multiplier $f$, minimize first in
$u$, and obtain
\begin{equation*}
\sup_ {\int f=0}\left\{\langle\eta_ h,f\rangle
-\frac12\int\nabla f^TH^{-1}\nabla f\right\}.
\end{equation*}
The maximizing potential solves $L_Hf=\eta_h$, and the value is
$\frac12\langle\eta_h,L_H^{-1}\eta_h\rangle$.  Since the covariance
of $\eta_h$ is $\Pi_h$, taking expectation proves the first identity
in (29).  The same computation with zero boundary values proves the
Dirichlet identity.

\smallskip\noindent
\underline{\textit{Step 2.}} Heat cutoff and the logarithmic coefficient.

It suffices to replace the cell projection in (29) by a heat cutoff at
time $n^{-2}$ with bounded error.  This exposes the logarithmic
contribution of the scales between $n^{-2}$ and $1$.
Write $P_n$ for projection onto the $n^2$ cell-constant functions and
$P_0$ for the constants, so $\Pi_h=P_n-P_0$.  Cell Poincare and
ellipticity give, for either boundary condition,
\begin{equation}
\|(I-P_ n)f\|_ 2^2\leq Cn^{-2}\|\nabla f\|_ 2^2
\leq C_ Hn^{-2}\langle L_ Hf,f\rangle.
\tag{32a}
\end{equation}
Consequently
\begin{equation}
\operatorname {Tr}((I-P_ n)e^{-tL_ H})
=\|(I-P_ n)e^{-tL_ H/2}\|_ {\rm HS}^2
\leq C_ Hn^{-2}\operatorname {Tr}(L_ He^{-tL_ H}).
\tag{32b}
\end{equation}
In the Neumann case all traces omit the constant mode.  The contribution
of $0<t<n^{-2}$ to
$\int_ 0^\infty\operatorname {Tr}(P_ ne^{-tL_ H})dt$ is at most one.
Integrating (32b) from $n^{-2}$ to infinity gives
\begin{equation}
C_ Hn^{-2}\int_ {n^{-2}}^\infty
\operatorname {Tr}(L_ He^{-tL_ H})dt
=C_ Hn^{-2}\operatorname {Tr}(e^{-n^{-2}L_ H})=O_ H(1).
\tag{32c}
\end{equation}
For Dirichlet conditions, replacing $P_n$ by $P_n-P_0$ subtracts the
bounded rank-one term
\[
|Q_1|^{-1}\langle1,(L_H^D)^{-1}1\rangle.
\]
The spectral gap absorbs the tail $t\geq1$.  Therefore
\begin{equation}
\operatorname {Tr}(\Pi_ hL_ H^{-1})
=\int_ {n^{-2}}^1\operatorname {Tr}'(e^{-tL_ H})dt+O_ H(1).
\tag{32d}
\end{equation}

The heat-trace coefficient $\sqrt{\det H}/(4\pi t)$, together with
the factor $1/2$ in the energy and the cutoff interval $[n^{-2},1]$,
gives
\begin{equation*}
\frac12\frac{\sqrt{\det H}}{4\pi}
\int_{n^{-2}}^1\frac{dt}{t}
=\frac{\sqrt{\det H}}{4\pi}\log n.
\end{equation*}
The factor $\sqrt{\det H}$ comes from the area of the frequency
ellipse for $L_H$.  We verify the heat-trace coefficient by the
following bracketing argument.  The linear change of variables turning
$L_H$ into the
Laplacian maps the square to a parallelogram.  Tile it by squares of side
$r$ and bracket by Dirichlet and Neumann direct sums over squares
contained in or meeting the parallelogram.  Interior lattice counts are
Riemann sums for
\begin{equation*}
(2\pi)^{-2}\int_ {\{\xi^TH^{-1}\xi\leq1\}}d\xi
=\frac{\sqrt{\det H}}{4\pi}.
\end{equation*}
The squares meeting the boundary contribute at most
$C_ H(rR+r^{-1}\sqrt R+r^{-2})$ eigenvalues below energy $R$.
Divide by $R$, first send $R\to\infty$, and then $r\downarrow0$.
Thus, for both boundary conditions,
\begin{equation}
N_ H(R)=\frac{\sqrt{\det H}}{4\pi}R+o_ H(R),
\tag{32e}
\end{equation}
uniformly on compact ellipticity classes.  Equivalently,
$t\operatorname {Tr}'(e^{-tL_ H})\to\sqrt{\det H}/(4\pi)$.
Logarithmic Cesaro averaging in (32d) yields
\begin{equation}
\frac12\operatorname {Tr}(\Pi_ hL_ H^{-1})
=\frac{\sqrt{\det H}}{4\pi}\log n+o_ H(\log n).
\tag{32}
\end{equation}
Since $\log n=h\log2$, (30) follows.

\smallskip\noindent
\underline{\textit{Step 3.}} Dirichlet--Neumann squeeze.

For the lower bound in (30b), paste the zero-boundary Dirichlet
potentials on the unit outer blocks.  Their common outer trace is
zero, so they are admissible in (30a), and their expected values add.
For the upper bound, paste the zero-normal Neumann flows instead.
Against any admissible full trace, integration by parts and quadratic
completion bound the dual value by the sum of their energies.  Both
bounds are proportional to the number $L^2$ of unit blocks, so the
normalization in (30a) makes the squeeze uniform in $L$.  No random
spatial mean is subtracted: the fields retain their zero-normal
boundary condition.

\smallskip\noindent
\underline{\textit{Step 4.}} Pointwise Gaussian moments.

The optimizers are linear images of the Gaussian source, so their
pointwise moments are determined by the sum of squared coefficients.
We will show that this sum is $O_H(1+\log n)$.  The spatial estimate
needed for both boundary conditions is the Green-gradient bound
$C_H|x-y|^{-1}$ in (32f): summing its squared cell contributions over
concentric annuli will produce the logarithm in (32h).

We first record the Green estimates and check that they apply after
the anisotropic change of variables.

\textbf{Dirichlet Green estimate.}  Diening, Rolfes and Salgado~\cite[Proposition 4.1]{DRS2024} state: if
$\Omega\subset\mathbb R^d$, $d\in\{2,3\}$, is a convex polytope and
$G:\Omega\times\Omega\to\mathbb R$ is its Green function (in the
paper's $W^{1,2}_ 0$ framework, hence the Dirichlet Green function of
the Laplacian), then for every $i$,
\begin{equation}
|\partial_ {x_ i}G(x,\xi)|\leq \frac{C_\Omega}{|x-\xi|^{d-1}}.
\tag{32fD}
\end{equation}
In addition, there is $\alpha\in(0,1]$, depending only on the inner
angles of $\Omega$, such that for every $i,j$ and all
$x,y,\xi\in\Omega$ off the displayed diagonals,
\begin{equation*}
\frac{|\partial_{x_i}G(x,\xi)-\partial_{y_i}G(y,\xi)|}{|x-y|^\alpha}
\leq C_ \Omega\bigl(
|x-\xi|^{-d-\alpha+1}+|y-\xi|^{-d+\alpha+1}\bigr),
\end{equation*}
\begin{equation*}
\frac{|\partial_{x_i}\partial_{\xi_j}G(x,\xi)-\partial_{y_i}\partial_{\xi_j}G(y,\xi)|}{|x-y|^\alpha}
\leq C_ \Omega\bigl(
|x-\xi|^{-d-\alpha}+|y-\xi|^{-d+\alpha}\bigr).
\end{equation*}
Only (32fD), the first assertion, is used below.  The proof of the cited
proposition points to the convex-polytope Green estimates identified there;
no smooth-boundary hypothesis is present.

\textbf{Neumann heat-gradient estimate.}  Wang and Yan~\cite[Theorem 1.2]{WangYan2010} use $P_ t=e^{t\Delta_ N}$ for the
Neumann heat semigroup on $\Omega$, $p_ t(x,y)$ for its kernel, and
$\|\cdot\|_ {1\to\infty}$ for the $L^1$-to-$L^\infty$ operator norm.
Its complete assertion is: if $\Omega$ is a compact convex domain in
$\mathbb R^d$, then there are $C,c>0$ such that, for every $t>0$,
\begin{equation}
\|\nabla P_ t\|_ {1\to\infty}\leq Ct^{-(d+1)/2},
\qquad
|\nabla_ xp_ t(x,y)|\leq Ct^{-(d+1)/2}
\exp\!\left(-\frac{|x-y|^2}{ct}\right).
\tag{32fN}
\end{equation}
The theorem explicitly covers nonsmooth convex domains.  Its proof
constructs reflecting Brownian motion on a convex domain, obtains the
gradient contraction and the short-time heat-kernel bound, and uses the
Neumann spectral gap for large time.  Section 4.3 of the same paper defines
the mean-zero Neumann Green operator and records the Green-gradient
consequence.  Indeed, in dimension two, with the normalization
\begin{equation*}
G^N_ \Omega(z,w)=\int_ 0^\infty
\{p_ t(z,w)-|\Omega|^{-1}\}\,dt,
\end{equation*}
the derivative of the constant vanishes and (32fN) is integrable for
$z\ne w$; differentiation first on a truncated time interval and then
dominated convergence give
\begin{equation}
|\nabla_ zG^N_ \Omega(z,w)|
\leq C\int_ 0^\infty t^{-3/2}
e^{-|z-w|^2/(ct)}\,dt
=\frac{C'}{|z-w|}.
\tag{32fN'}
\end{equation}

We now check the hypotheses and the change of variables.  Let
$S=H^{1/2}$ and $\mathcal P_ H=SQ_ 1$.  The variational change
$z=Sx$ transforms
$L_ H=-\operatorname {div}_ x(H^{-1}\nabla_ x)$ into the Euclidean
Laplacian on the parallelogram $\mathcal P_ H$, preserving respectively
Dirichlet and natural Neumann boundary conditions.  The parallelogram is
a compact convex polygon, so (32fD) and (32fN') apply, including at its
corners.  Green-kernel normalization gives, for either boundary condition,
\begin{equation*}
G_ H^\#(x,y)=(\det S)G_ {\mathcal P_ H}^\#(Sx,Sy).
\end{equation*}
Therefore
\begin{equation}
|H^{-1}\nabla_ xG_ H^\#(x,y)|
=\det(S)|H^{-1/2}\nabla_ zG_ {\mathcal P_ H}^\#(Sx,Sy)|
\leq C_ H|x-y|^{-1}.
\tag{32f}
\end{equation}
Thus (32f) holds for the Dirichlet kernel and the mean-zero Neumann
kernel.  Its constant is allowed to depend on the fixed
positive-definite matrix $H$, as in (31).  We now apply this decay
bound to the cell coefficients.

If $e_ C=n\mathbf1_ C$, the coefficient of the cell Gaussian $Z_ C$
in the primal optimizer is
\begin{equation}
k_ C(x)=n\left(\int_ C H^{-1}\nabla_ xG_ H^\#(x,y)\,dy
-n^{-2}\int_ {Q_ 1}H^{-1}\nabla_ xG_ H^\#(x,y)\,dy\right).
\tag{32g}
\end{equation}
For every cell meeting $B(x,2/n)$, polar integration gives
$\int_ C|x-y|^{-1}dy\leq C/n$; hence its first term in (32g) is
$O_ H(1)$, and the boundedly many such cells contribute $O_ H(1)$ to
$\sum_ C|k_ C(x)|^2$.  In the annulus
$j/n\leq\operatorname {dist}(x,C)<(j+1)/n$, there are $O(j)$
cells and (32f) gives
$n\int_ C|x-y|^{-1}dy\leq C/j$; together with the $O(1/n)$
mean term estimated below, $|k_ C(x)|\leq C_ H/j$.  Summing
$\sum_ {j\leq n}j\cdot j^{-2}$ gives
\begin{equation}
\sup_ x\sum_ C|k_ C(x)|^2\leq C_ H(1+\log n).
\tag{32h}
\end{equation}
For the mean-subtraction term,
$n^{-1}\int_ {Q_ 1}|x-y|^{-1}dy\leq C/n$ uniformly in $x$, so its
squared contribution summed over all $n^2$ coefficients is $O(1)$.
Gaussian moment identities, followed by spatial integration, prove
(31) for both primal and dual optimizers.

\end{proof}

\subsection{Upper consistency from a Neumann flux}

The upper estimate uses the profile only near the point being tested.
A quadratic Neumann optimizer has size of order $\sqrt h$, whereas
the neighborhood of the expansion has radius of order $\sqrt T$.
Since $h/T\to0$, the Gaussian tail outside this neighborhood is
negligible.  The argument applies also when the contact slope has
norm one.

\begin{lemma}\label{lem:localized-upper-consistency}
Put
\begin{equation}
\mathcal G(H)=\frac1{4\pi}\sqrt{\det H},\qquad H\in\mathbb S_ +^2.
\tag{33}
\end{equation}
Let $F_ T$ be a convex $1$-Lipschitz profile, $p_ T=\sqrt Tq_ 0$,
and suppose that, for $|z|\leq r_ 0\sqrt T$,
\begin{equation}
F_ T(p_ T+z)-F_ T(p_ T)
\leq a\cdot z+\frac1{2\sqrt T}z^THz+C\frac{|z|^3}{T}.
\tag{33a}
\end{equation}
There is no restriction on $|a|\leq1$.  If $H>0$, $h\to\infty$,
and $h/T\to0$, then, uniformly in $L$,
\begin{equation}
\begin{split}
(\mathcal H_ {h,L}F_ T)(p_ T)-F_ T(p_ T)
\leq{}&\frac{h\log2}{\sqrt T}\mathcal G(H)+\frac{o_ H(h)}{\sqrt T}\\
&+C_ H\frac{h^{3/2}}T+C_ H\sqrt h\,e^{-c_ HT/h}.                  
\end{split}
\tag{33b}
\end{equation}
For singular $H$, the same upper consistency follows by using
$H+\gamma I$, taking the scaling limit first, and then
$\gamma\downarrow0$.

\end{lemma}

\begin{proof}
For every realization and every zero-normal flux $u$ with
$-\operatorname {div}u=\eta_ h^{(L)}$, integration by parts and
Fenchel's inequality give the general primal domination
\begin{equation}
(\mathcal H_ {h,L}F_ T)(p)
\leq\mathbb E\fint_ {Q_ L}F_ T(p+u(x))\,dx.
\tag{33c}
\end{equation}
Choose on every unit outer block the Neumann quadratic minimizer for
$\frac12\int u^THu$, and paste these zero-normal fields.  This is a
linear centred Gaussian function of the source; hence
$\mathbb E\fint u=0$.  Let
$B=\{(\omega,x):|u(\omega,x)|>r_ 0\sqrt T\}$.  On $B^c$, insert
(33a); on $B$, use the $1$-Lipschitz bound.  The full linear term
vanishes in expectation, so its contribution on $B^c$ can be estimated
on $B$ through the identity
\begin{equation*}
\mathbb E\fint 1_ {B^c}a\cdot u
=-\mathbb E\fint1_ Ba\cdot u.
\end{equation*}
This cancellation uses the whole probability and spatial domain; the
linear term need not vanish on $B^c$ alone.

The local upper inequality in both signs forces $|a|\leq1$.
Hence the transferred linear term and the Lipschitz contribution on
$B$ are both bounded by $\mathbb E\fint1_ B|u|$.  The positive
quadratic term on $B^c$ is at most its full-space value, evaluated by
(30), while (31) bounds the cubic remainder by $C_ Hh^{3/2}/T$.
At every fixed $x$, (32h) gives
$\operatorname {Var}u(x)\leq C_ Hh$.  Gaussian tails, together with
the preceding two bad-event bounds, give
$C_ H\sqrt h\,e^{-c_ HT/h}$.  This proves (33b) for $H>0$.  If $H$ is singular,
the Taylor quadratic is bounded above by the one with $H+\gamma I$;
apply the preceding estimate with fixed $\gamma>0$, use (30), and only
then send $\gamma\downarrow0$.

\end{proof}

\subsection{Consistency from a truncated Dirichlet potential}

For a lower bound we need an admissible dual potential, not just a
bound on its expected gradient.  The profiles below have a conjugate
cost that is exactly quadratic near the contact slope and lies above
that quadratic everywhere.  The global inequality bounds the cell
value from above.  To obtain the matching lower bound, we truncate a
Dirichlet optimizer so that every gradient lies in the region of
quadratic equality.  These are the profiles used for the lower
barriers in the next subsection.

\begin{lemma}\label{lem:nonlinear-consistency}
Use $\mathcal G$ from (33).  Let $\psi$ be either a global convex
$1$-Lipschitz cone-growth profile or a proper closed convex profile
whose conjugate domain is compactly contained in $B_ 1$,
let $q_ 0\in\mathbb R^2$, and suppose
$a_ 0\in\partial\psi(q_ 0)$ lies in the open unit ball.  Assume that,
on $\overline B_ 1-a_ 0$, its conjugate has the exact local form
\begin{equation}
\psi^*(a_ 0+b)=\psi^*(a_ 0)+q_ 0\cdot b+D(b),
\tag{34}
\end{equation}
where $D(0)=0$, $D$ is convex, and, for some $\rho>0$,
\begin{equation}
D(b)=\tfrac12b^TH^{-1}b\quad\text{if }|b|_ {H^{-1}}\leq\rho.
\tag{34a}
\end{equation}
Assume on its whole effective domain that
\begin{equation*}
D(b)\geq Q(b):=\tfrac12b^TH^{-1}b;
\end{equation*}
thus $D$ is the quadratic near zero plus a convex exterior wall.  Put
$\Psi_ T(p)=\sqrt T\,\psi(p/\sqrt T)$ and
$p_ T=\sqrt Tq_ 0$.  If $H>0$, $h\to\infty$, and
$T/h\to\infty$, then, uniformly in $L$,
\begin{equation}
\begin{split}
&(\mathcal H_ {h,L}\Psi_ T)(p_ T)-\Psi_ T(p_ T)\\
&\quad=\frac{h\log2}{\sqrt T}\,\mathcal G(H)+\frac{o_ H(h)}{\sqrt T}
+O_ H\!\left(\frac{h^2}{T^{3/2}}\right).
\end{split}
\tag{35}
\end{equation}
This lemma is used only with $H>0$ and an interior dual slope.

\end{lemma}

\begin{proof}
After removing the affine part in (34), we will bound the increment
between the Dirichlet and Neumann quadratic values, divided by
$\sqrt T$, with an additional lower-bound error of order
$h^2/T^{3/2}$.  Since both quadratic values have the leading term
(30), this proves (35).  The only change needed in the quadratic
Dirichlet competitor is a truncation that keeps its gradient in the
set where $D=Q$.

\smallskip\noindent
\underline{\textit{Step 1.}} Dual reduction and quadratic value.

Since $\Psi_ T^*=\sqrt T\,\psi^*$, the definition (16) and (34) turn
the increment on the left of (35) exactly into
\begin{equation}
\mathcal V_ D=L^{-2}\max_ c\mathbb E\sup_ {f|_ {\partial Q_ L}=c}
\left\{\int f\eta_ h^{(L)}-\sqrt T\int D(\nabla f)\right\}.
\tag{36}
\end{equation}
The term $p_ T\cdot a$ cancels the spatial average of the linear term
in (34), and the remaining affine source term has expectation zero.  For
$Q(b)=\frac12b^TH^{-1}b$, the change of variable
$g=\sqrt T f$ and (30a) give
\begin{equation}
\mathcal V_ Q=T^{-1/2}\mathcal E^{\rm tr}_ {h,L}(H).
\tag{37}
\end{equation}
Lemma~\ref{lem:grouped-quadratic-generator} evaluates this quadratic
value.  To obtain a lower bound for the nonlinear value, we now keep
the Dirichlet competitor inside the region where the costs agree.

\smallskip\noindent
\underline{\textit{Step 2.}} Truncated Dirichlet competitor.

Let $f_D$ be the zero-boundary Dirichlet dual optimizer for the
unscaled quadratic problem, pasted over the unit blocks, and put
$b_D=T^{-1/2}\nabla f_D$.  Equation (31) gives
\begin{equation}
\mathbb E\fint|b_ D|^4\leq C_ H(h/T)^2,\qquad
\sup_ x\mathbb E|b_ D(x)|^2\leq C_ Hh/T.
\tag{37a}
\end{equation}
These moment bounds do not ensure that the Gaussian optimizer stays
pointwise in the set where $D=Q$.  Apply
Lemma~\ref{lem:lipschitz-truncation} to $f_ D/\sqrt T$ at a threshold
$\lambda$, fixed independently of $h$ and $T$, so that its resulting
bound $C\lambda$ is
smaller than both
$\rho/\|H^{-1/2}\|$ and
$\operatorname {dist}(a_ 0,\partial B_ 1)$.  The zero extension is
globally Lipschitz and belongs to $W^{1,4}_ 0$, so the truncated block
potentials paste with exactly zero trace.  By (28) and (37a), the
expected squared gradient error is at most $C_H(h/T)^2$ per unit
area.  Quadratic completion multiplies this error by $\sqrt T$,
the coefficient of the cost in (36).  Thus the lost value is at most
$C_Hh^2/T^{3/2}$.

\smallskip\noindent
\underline{\textit{Step 3.}} Comparison of the two values.

Since $D\geq Q$, (37) and the upper bound in (30b) bound
$\mathcal V_D$ by $T^{-1/2}E_{h,H}^N$.  The truncated competitor has
zero trace and stays everywhere in the region where $D=Q$, so its
value is at least $T^{-1/2}E_{h,H}^D$ minus the preceding error.
Applying (30) to these two bounds proves (35).  The construction and
the error estimate are blockwise, hence uniform in $L$.

\end{proof}

\subsection{Global barriers for local lower contacts}

Monotonicity of the cell operator compares profiles on the whole
space, whereas a viscosity test initially gives only a local order
relation.  We therefore replace a local lower test by a global convex
profile below the spatial slice.  It must retain almost the same
Hessian at the contact, meet the hypotheses of
Lemma~\ref{lem:nonlinear-consistency}, and grow strictly more slowly
than the slice at infinity.  The last property keeps the later
minimizing points in a compact set.  The uniform estimate (38c) is
needed because these points will approach, rather than equal, the
original contact point.

\begin{lemma}\label{lem:global-barriers}
Let $U$ be convex and $1$-Lipschitz, with
$|p|-C\leq U(p)\leq|p|+C$.  Suppose a convex $C^2$ function
$\phi$ touches $U$ strictly from below at $q_ 0$,
$|a_ 0|:=|\nabla\phi(q_ 0)|<1$, and
$H=D^2\phi(q_ 0)>0$.  For every
$0<\delta<\lambda_ {\min}(H)$ there is a global proper convex barrier
$B_ \delta\leq U$, touching at $q_ 0$, such that
\begin{equation}
B_ \delta(q)<U(q)\quad(q\ne q_ 0),\qquad
D^2B_ \delta(q_ 0)=H-\delta I,
\tag{38a}
\end{equation}
and
\begin{equation}
\operatorname {dom}B_ \delta^*\Subset B_ 1,\qquad
U(q)-B_ \delta(q)\longrightarrow+\infty
\quad\text{as }|q|\to\infty.
\tag{38b}
\end{equation}
Near $a_ 0$, its conjugate is exactly quadratic with Hessian
$(H-\delta I)^{-1}$.  Its conjugate Bregman gap---that is,
$B_\delta^*(a+b)-B_\delta^*(a)-\nabla B_\delta^*(a)\cdot b$---based at every
slope sufficiently near $a_ 0$, is bounded below globally by this
same quadratic form and equals it for small increments.  Consequently,
for
$(B_ \delta)_ T(p)=\sqrt T B_ \delta(p/\sqrt T)$,
\begin{equation}
(\mathcal H_ {h,L}(B_ \delta)_ T)(\sqrt Tq_ 0)
-(B_ \delta)_ T(\sqrt Tq_ 0)
=\frac{h\log2}{\sqrt T}\mathcal G(H-\delta I)+\frac{o_ {H,\delta}(h)}{\sqrt T}
+O_ {H,\delta}\!\left(\frac{h^2}{T^{3/2}}\right).
\tag{38}
\end{equation}
Moreover, there is a neighborhood $V$ of $q_ 0$ on which
\begin{equation}
\sup_ {q\in V}\left|
\frac{\sqrt T}{h\log2}
\left\{(\mathcal H_ {h,L}(B_ \delta)_ T)(\sqrt Tq)
-(B_ \delta)_ T(\sqrt Tq)\right\}
-\mathcal G(H-\delta I)\right|\longrightarrow0
\tag{38c}
\end{equation}
whenever $h\to\infty$ and $h/T\to0$, uniformly in $L$.
No global upper barrier is asserted or needed.

\end{lemma}

\begin{proof}
Write $g=U^*$.  Since conjugation reverses order, it is enough to
construct a convex function above $g$ whose conjugate has the desired
contact.  The local test gives a quadratic upper bound for $g$ near
$a_0$.  We extend that bound over a compact ball inside $B_1$, with
strict inequality away from $a_0$, and set it equal to $+\infty$
outside the ball.  This choice will give both the global order and
the slower growth required in (38b).

\smallskip\noindent
\underline{\textit{Step 1.}} The dual contact and a strict quadratic bound.

The conjugate $g$ is finite and continuous in $B_ 1$ and is
$+\infty$ off $\overline B_ 1$.  The local lower contact and
convexity give $a_ 0\in\partial U(q_ 0)$, so $q_ 0$ maximizes
\begin{equation}
g(a_ 0)=\sup_ q\{a_ 0\cdot q-U(q)\}.
\tag{39}
\end{equation}
It is the unique maximizer.  Indeed, a second maximizer would make
$U$ affine with slope $a_ 0$ on the segment joining it to $q_ 0$.
The positive-definite quadratic term of the lower test would then put
$\phi$ strictly above that affine function immediately away from
$q_ 0$, contrary to the local lower contact.

Choose $r<1$ so that a neighborhood of $a_ 0$ is contained in
$B_ r$.  From the lower cone bound,
\begin{equation}
a\cdot q-U(q)\leq-(1-r)|q|+C\qquad(a\in B_ r).
\tag{40}
\end{equation}
Thus all maximizers defining $g(a)$ for $a$ near $a_ 0$ lie in
one compact set.  Uniqueness in (39) further implies that these
maximizers converge to $q_ 0$ as $a\to a_ 0$.  They consequently
lie, after the slope neighborhood is reduced, in the neighborhood
where $U\geq\phi$.  The local Legendre transform of $\phi$ has
Hessian $H^{-1}$ at $a_ 0$.  Comparing first with
$H-\frac\delta2I$, and only then with $H-\delta I$, gives, for
all sufficiently small $v$,
\begin{equation}
\begin{split}
g(a_ 0+v)-g(a_ 0)-q_ 0\cdot v
&\leq\frac12v^T(H-\tfrac\delta2I)^{-1}v\\
&\leq\frac12v^T(H-\delta I)^{-1}v-c_ \delta|v|^2
\end{split}
\tag{41}
\end{equation}
for some $c_ \delta>0$.  The strict second inequality is the
ingredient that prevents any further contact.

\smallskip\noindent
\underline{\textit{Step 2.}} Extension beyond the contact slope.

We now extend the quadratic bound while leaving it unchanged near
$a_0$.  Choose the radius of the quadratic region and the compact
dual domain first; then choose the coefficient of a convex term to
dominate the remaining annulus.

Put $A_ \delta=(H-\delta I)^{-1}$ and
$Q_ \delta(v)=\frac12v^TA_ \delta v$.  Choose
\begin{equation*}
0<2\rho<d<\operatorname {dist}(a_ 0,\partial B_ 1)
\end{equation*}
so that (41) holds for $|v|\leq2\rho$, and, on $|v|\leq d$, set
\begin{equation}
D(v)=Q_ \delta(v)+C_ \rho(|v|-\rho)_ +^2.
\tag{41a}
\end{equation}
This is convex and equals $Q_ \delta$ on $|v|\leq\rho$.  Choose
$C_ \rho>0$ strictly larger than
\begin{equation}
\max\!\left\{0,\
\sup_ {2\rho\leq|v|\leq d}
\frac{g(a_ 0+v)-g(a_ 0)-q_ 0\cdot v-Q_ \delta(v)}
{(|v|-\rho)^2}\right\}.
\tag{41b}
\end{equation}
The supremum is finite by continuity on the compact annulus.
Equations (41)--(41b) imply
\begin{equation}
D(v)>g(a_ 0+v)-g(a_ 0)-q_ 0\cdot v
\qquad(0<|v|\leq d).
\tag{41c}
\end{equation}
\smallskip\noindent
\underline{\textit{Step 3.}} The global lower barrier.

Define the proper closed convex function
\begin{equation}
\widetilde g(a_ 0+v)=g(a_ 0)+q_ 0\cdot v+D(v)
\quad(|v|\leq d),\qquad
\widetilde g=+\infty\quad(|v|>d),
\tag{41d}
\end{equation}
and set $B_ \delta=\widetilde g^*$.  We have
$\widetilde g\geq g$, with equality among its effective slopes only
at $a_ 0$; hence $B_ \delta\leq U$.  At $q_ 0$, the expression
defining $B_ \delta(q_ 0)$ is $U(q_ 0)-D(v)$, so its unique maximizing
slope is $a_ 0$ and $B_ \delta(q_ 0)=U(q_ 0)$.

If equality $B_ \delta(q)=U(q)$ held elsewhere, a maximizing slope
for the compact-domain conjugate $B_ \delta(q)$ would have to be a
slope where $\widetilde g=g$, hence $a_ 0$.  But
$\widetilde g$ is differentiable at $a_ 0$, with gradient $q_ 0$,
so $a_ 0$ can maximize its conjugate only at $q=q_ 0$.  This proves
the strict assertion in (38a).  Since (41d) is exactly affine plus
$Q_ \delta$ near $a_ 0$, local conjugacy gives
$D^2B_ \delta(q_ 0)=H-\delta I$.  Also
$\operatorname {dom}B_ \delta^*=\overline B_ d(a_ 0)\Subset B_ 1$,
so $B_ \delta(q)\leq r_ \delta|q|+C_ \delta$ for some
$r_ \delta<1$.  The lower cone bound for $U$ proves (38b).

\smallskip\noindent
\underline{\textit{Step 4.}} Uniform Bregman gap and consistency.

The discrete contact will occur at slopes near $a_0$, so consistency
at $a_0$ alone is insufficient.  We verify the same quadratic-gap
hypotheses uniformly at these nearby slopes.  Write the wall in
(41a) as $W$.  If $|a-a_0|<\rho/4$, then, for every increment
$b$ for which the left side is finite,
\begin{equation}
\begin{split}
&B_ \delta^*(a+b)-B_ \delta^*(a)
-\nabla B_ \delta^*(a)\cdot b\\
&\qquad=Q_ \delta(b)+
\{W(a-a_ 0+b)-W(a-a_ 0)-\nabla W(a-a_ 0)\cdot b\}
\geq Q_ \delta(b).                                             
\end{split}
\tag{41e}
\end{equation}
For increments outside the compact conjugate domain the left side is
$+\infty$.  Equality holds for all sufficiently small $b$,
uniformly for $a$ in a smaller ball.  Thus the hypotheses and all
constants in Lemma~\ref{lem:nonlinear-consistency} hold uniformly for the
corresponding spatial points $q$ near $q_ 0$.  The strict-slope
extension in Definition~\ref{def:cell-operator} has an attained cell
value and is monotone, as proved in Proposition~\ref{prop:cell-operator-properties}.  Lemma~\ref{lem:nonlinear-consistency} gives
(38), and, after division by $h\log2/\sqrt T$, gives (38c), because
its two normalized errors are $o(1)$ and $O(h/T)$.

\end{proof}

\section{The convex viscosity limit}\label{sec:viscosity}

We now show that many small cell-operator increments converge to the
evolution with generator $\mathcal G$.  There are two tasks: identify
the equation satisfied in the limit, and prove uniqueness in the class
of convex profiles with bounded corrections to $|q|$.  The first uses
the local Gaussian upper estimate and the global lower barriers from
the preceding section.  The second requires a comparison argument on
$\mathbb R^2$, where a maximizing pair of spatial points can escape
to infinity together.

\begin{lemma}\label{lem:viscosity-limit}
Let $h_ m\to\infty$, $r_ m\to\infty$, and $h_ m/r_ m\to0$, and let
$L_ m$ be arbitrary integers tending to infinity.  For each $m$, define the
$m$th row of the triangular monotone scheme by
\begin{equation}
F_0^{(m)}(p)=|p|,\qquad
F_{k+1}^{(m)}=\mathcal H_{h_m,L_m}F_k^{(m)},
\qquad t_k^{(m)}=k h_m\log2.
\tag{42}
\end{equation}
Thus $h_m$ and $L_m$ are fixed during the iteration in row $m$, while the
row parameters vary with $m$.  Use piecewise-constant interpolation in the
physical time $t$ and then rescale time by $r_m$; equivalently, $U_m$ below is
constant on each interval
$[t_k^{(m)}/r_m,t_{k+1}^{(m)}/r_m)$.  Thus
\begin{equation}
U_ m(s,q)=r_ m^{-1/2}F_{\lfloor sr_ m/(h_ m\log2)\rfloor}^{(m)}
(\sqrt{r_ m}\,q)
\tag{43}
\end{equation}
converges locally uniformly, at every fixed positive time, to the unique
convex cone-growth viscosity solution of
\begin{equation}
\partial_ sU=\mathcal G(D_ q^2U)=\frac1{4\pi}\sqrt{\det D_ q^2U},
\qquad U(0,q)=|q|.
\tag{44}
\end{equation}
Here the viscosity formulation is the strong convex one:
$U(s,\cdot)$ is convex; every upper test, and every lower test whose
spatial Hessian is positive semidefinite, satisfies the corresponding
evolution inequality, including when that Hessian is singular.  In
particular, with $\kappa=U(1,0)$, every integer sequence $k_m$ satisfying
the conditions below obeys
\begin{equation}
\frac{F_{k_m}^{(m)}(0)}{\sqrt{t_{k_m}^{(m)}}}\longrightarrow\kappa
\quad\hbox{whenever }t_{k_m}^{(m)}\to\infty,
\quad h_m/t_{k_m}^{(m)}\to0.
\tag{45}
\end{equation}

\end{lemma}

\begin{proof}
We first obtain the time modulus and identify the two viscosity
inequalities.  For an upper contact the Gaussian flux uses only the
local test.  For a lower contact we minimize the difference from a
global spatial barrier; the resulting scalar function of time retains
the contact derivative and satisfies the discrete inequality (47i).
We then prove comparison.  Translating a maximizing pair by its
common center allows us to remove the spatial localization while
retaining the matrix order supplied by the relative quadratic
penalty.  Uniqueness and scaling will give (45).

\smallskip\noindent
\underline{\textit{Step 1.}} Stability, compactness, and the initial trace.

The one-step bound (16a), summed directly, is too large for the time
modulus we need.  We instead iterate its primal bound using
independent centered fluxes, whose second moments add.
Let $Q$ be the pasted linear fine flux in (16b), let $X$ be uniform
on $Q_{L_m}$, independent of the source, and put $Z=Q(X)$.  Then
$\mathbb EZ=0$ and
$\mathbb E|Z|^2\leq Ch_ m$.  Equation (16b), before using the
Lipschitz estimate, gives the stronger domination
\begin{equation}
(\mathcal H_ {h_ m,L_ m}\Psi)(p)\leq\mathbb E\Psi(p+Z)
\tag{45a}
\end{equation}
for every admissible profile $\Psi$.  Use independent copies of the
fine source, the uniform point, and hence $Z$, at successive
iterations.  Repeated application of (45a) within row $m$ gives
\begin{equation*}
F_{k+\ell}^{(m)}(p)
\leq\mathbb E F_k^{(m)}(p+Z_ 1+\cdots+Z_ \ell).
\end{equation*}
Since every $F_k^{(m)}$ is $1$-Lipschitz,
\begin{equation}
0\leq F_{k+\ell}^{(m)}(p)-F_k^{(m)}(p)\leq C\sqrt{\ell h_m},
\qquad |p|\leq F_k^{(m)}(p)\leq|p|+C\sqrt{k h_m}.
\tag{46}
\end{equation}
Consequently, at grid times,
\begin{equation}
|U_ m(s,q)-U_ m(t,q)|
\leq C\sqrt{|s-t|}+C\sqrt{h_ m/r_ m},
\tag{46a}
\end{equation}
while the spatial Lipschitz constant is one.  Thus (43) is locally
equicontinuous and
$0\leq U_ m(s,q)-|q|\leq C\sqrt s+o(1)$, which gives the initial
trace uniformly on $\mathbb R^2$.

\smallskip\noindent
\underline{\textit{Step 2.}} Upper half-relaxed limit and the subsolution inequality.

For consistency put $\Delta_m=h_m\log2/r_m$ and
$U_m^k(q):=U_m(k\Delta_m,q)$, and define
\begin{equation}
(S_ m\varphi)(q)=r_ m^{-1/2}\mathcal H_ {h_ m,L_ m}
[p\mapsto\sqrt{r_ m}\,\varphi(p/\sqrt{r_ m})](\sqrt{r_ m}q).
\tag{46b}
\end{equation}
The map is monotone and constant-equivariant.
We use the standard half-relaxed limits
\[
\overline U(s,q)=\limsup_{\substack{m\to\infty\\(s',q')\to(s,q)}}
U_m(s',q'),\qquad
\underline U(s,q)=\liminf_{\substack{m\to\infty\\(s',q')\to(s,q)}}
U_m(s',q').
\]

Let a smooth convex $\phi$ touch the upper half-relaxed limit strictly
from above at $(s_0,q_0)$, $s_0>0$.  We must bound its time derivative
by $\mathcal G(D_q^2\phi)$ using only this local contact.
Select grid maxima
$(s_ m,q_ m)=(k_ m\Delta_ m,q_ m)\to(s_ 0,q_ 0)$ in a fixed small
parabolic cylinder.  If $V_ m$ is the pasted Neumann quadratic flux for
$D_ q^2\phi(s_ 0,q_ 0)+\gamma I$, the primal domination (33c), expressed
in the rescaled variables, gives
\begin{equation}
U_ m^{k_ m}(q_ m)
\leq\mathbb E\fint
U_ m^{k_ m-1}(q_ m+V_ m(x)/\sqrt{r_ m})\,dx.
\tag{46c}
\end{equation}
On $|V_ m|\leq r_ 0\sqrt{r_ m}$, maximality in the cylinder bounds the
right side by the same expectation with $U_ m$ replaced by
$\phi$, plus the contact constant.  On the complementary event, the
cone bounds (46), the spatial Lipschitz bound, and (32h) give
$C\sqrt{h_ m/r_ m}\exp(-c r_ m/h_ m)$, which is
$o(\Delta_ m)$.  Taylor expansion of
\begin{equation*}
\sqrt{r_ m}\{\phi(s_ m-\Delta_ m,q_ m+z/\sqrt{r_ m})
-\phi(s_ m-\Delta_ m,q_ m)\}
\end{equation*}
has exactly the form (33a).  The centred linear term vanishes.  Lemma~\ref{lem:localized-upper-consistency} therefore yields
\begin{equation}
\frac{\phi(s_ m,q_ m)-\phi(s_ m-\Delta_ m,q_ m)}{\Delta_ m}
\leq \mathcal G(D_ q^2\phi(s_ 0,q_ 0)+\gamma I)+o(1).
\tag{46d}
\end{equation}
First let $m\to\infty$ and then $\gamma\downarrow0$.  This proves
\begin{equation}
\phi_ s(s_ 0,q_ 0)\leq \mathcal G(D^2\phi(s_ 0,q_ 0)),
\tag{47}
\end{equation}
including unit slopes and singular upper jets.  No global upper barrier
is used.

\smallskip\noindent
\underline{\textit{Step 3.}} Lower contact and the scalar minimum gap.

For a lower contact, the test itself need not lie below the preceding
time slice on the whole space, so it cannot be inserted directly into
the monotone scheme.  We use Lemma~\ref{lem:global-barriers} to fix
a global spatial barrier and allow only its additive constant to vary
with time.  For every $S<\infty$, the stability estimates pass to the
limit on $0\leq s,t\leq S$:
\begin{equation}
|\underline U(s,q)-\underline U(t,q)|
\leq C\sqrt{|s-t|},\qquad
|q|\leq\underline U(s,q)\leq|q|+C_ S.
\tag{47a}
\end{equation}
Let a smooth convex $\phi$ touch $\underline U$ strictly from
below at $(s_ 0,q_ 0)$, with $s_ 0>0$, and put
\begin{equation}
\tau=\phi_ s(s_ 0,q_ 0),\qquad
a_ 0=D_ q\phi(s_ 0,q_ 0),\qquad
H=D_ q^2\phi(s_ 0,q_ 0)>0.
\tag{47b}
\end{equation}
First suppose $|a_ 0|<1$.  Fix
$0<\delta<\lambda_ {\min}(H)$ and apply Lemma~\ref{lem:global-barriers} to the spatial slice
$q\mapsto\underline U(s_ 0,q)$, using
$q\mapsto\phi(s_ 0,q)$ as its lower test.  Denote the resulting
barrier by $B_ \delta$.

The largest constant that can be added to this barrier at time $s$ is
\begin{equation}
c(s)=\min_ {q\in\mathbb R^2}
\{\underline U(s,q)-B_ \delta(q)\}.
\tag{47c}
\end{equation}
The minimum is attained, uniformly for $s$ near $s_ 0$, because
(47a) and (38b) make the expression uniformly coercive.  Strict global
contact gives
\begin{equation}
c(s_ 0)=0,\qquad
\operatorname {argmin}_ q
\{\underline U(s_ 0,q)-B_ \delta(q)\}=\{q_ 0\}.
\tag{47d}
\end{equation}
The common time modulus, strictness on compact annuli, and coercivity
then imply that every minimizer in (47c) converges to $q_ 0$ as
$s\to s_ 0$.

To preserve the original time derivative $\tau$, we need a smooth
lower test for $c$ with that derivative at $s_0$.  More precisely,
there are $K,r>0$ such that
\begin{equation}
c(s)\geq\ell(s):=\tau(s-s_ 0)-K(s-s_ 0)^2
\quad\text{for }|s-s_ 0|\leq r,\qquad c(s_ 0)=\ell(s_ 0).
\tag{47e}
\end{equation}
Indeed, all relevant minimizers lie in the parabolic neighborhood
where $\underline U\geq\phi$.  Put $z=q-q_ 0$ and
$\sigma=s-s_ 0$.  The exact local quadratic formula for $B_ \delta$
and the spacetime Taylor expansion of $\phi$ give
\begin{equation}
\phi(s,q)-B_ \delta(q)
=\tau\sigma+\frac\delta2|z|^2
+\sigma D^2_ {sq}\phi(s_ 0,q_ 0)\cdot z
+O(\sigma^2)+o(|z|^2).
\tag{47f}
\end{equation}
After shrinking $r$, absorb the last spatial remainder into
$\delta|z|^2/4$.  Completing the square absorbs the mixed term at a
cost of order $\sigma^2$, which proves (47e).  The small Hessian loss
$\delta I$ in the barrier is what permits this absorption.

\smallskip\noindent
\underline{\textit{Step 4.}} Discrete contacts and the one-step inequality.

It suffices to compare the time increment of this minimum with the
one-step increment of $B_\delta$.  Temporal contact will bound the
former from above by the derivative of $\ell$; global spatial order
and monotonicity will bound it from below by the latter.  The uniform
consistency (38c) then supplies the desired inequality for $\tau$.

At grid times put
\begin{equation}
c_ m^k=\min_ q\{U_ m^k(q)-B_ \delta(q)\},
\tag{47g}
\end{equation}
and choose a minimizer $q_ m^k$.  These minima are attained, and their
minimizers are uniformly bounded near $s_ 0$, since
\begin{equation*}
U_ m^k(q)-B_ \delta(q)
\geq(1-r_ \delta)|q|-C_ \delta.
\end{equation*}
The lower half-relaxed limit of the scalar grid functions $c_ m$ is
exactly $c$.  For the lower bound, take minimizing points, use their
compactness, and invoke the definition of $\underline U$.  For the
reverse bound, take a minimizer of (47c), use a recovery sequence for
$\underline U$ there, and move its time coordinate to the nearest
grid time using (46a).

Fix $\nu>0$ and replace $\ell$ on
$[s_ 0-r,s_ 0+r]$ by
$\ell_ \nu(s)=\ell(s)-\nu(s-s_ 0)^4$.  Then
$c-\ell_ \nu$ has its unique minimum at $s_ 0$, and
$\ell_ \nu'(s_ 0)=\tau$.  Minimizing
$c_ m^k-\ell_ \nu(k\Delta_ m)$ over the grid points in this interval
and using the preceding half-relaxed-limit equality produces interior
indices $k_ m$ such that
\begin{equation}
t_ m=k_ m\Delta_ m\longrightarrow s_ 0,\qquad
c_ m^{k_ m}-\ell_ \nu(t_ m)
\leq c_ m^j-\ell_ \nu(j\Delta_ m)
\tag{47h}
\end{equation}
for every grid point $j\Delta_ m$ in the interval.  If
$q_ m=q_ m^{k_ m}$, every subsequential limit of $q_ m$, together with
the recovery upper bound implicit in this selection, attains the
minimum in (47d).  Hence $q_ m\to q_ 0$.

By the definition of $c_ m^{k_ m-1}$,
\begin{equation*}
U_ m^{k_ m-1}\geq B_ \delta+c_ m^{k_ m-1}.
\end{equation*}
Using $U_ m^{k_ m}=S_ mU_ m^{k_ m-1}$, monotonicity and constant
equivariance, and then evaluating at the current minimizer $q_ m$,
gives the exact recursion
\begin{equation}
c_ m^{k_ m}-c_ m^{k_ m-1}
\geq S_ mB_ \delta(q_ m)-B_ \delta(q_ m).
\tag{47i}
\end{equation}
On the other hand, (47h), applied also at $k_ m-1$, gives
\begin{equation}
c_ m^{k_ m}-c_ m^{k_ m-1}
\leq\ell_ \nu(t_ m)-\ell_ \nu(t_ m-\Delta_ m).
\tag{47j}
\end{equation}
Divide by $\Delta_ m$.  Since $q_ m\to q_ 0$, the uniform consistency
(38c), written in the variables (46b), makes the right side of (47i)
converge to $\mathcal G(H-\delta I)$; the quotient in (47j) converges to
$\ell_ \nu'(s_ 0)=\tau$.  Therefore
\begin{equation*}
\phi_ s(s_ 0,q_ 0)\geq \mathcal G(H-\delta I).
\end{equation*}
Letting $\delta\downarrow0$ proves the supersolution inequality for
positive-definite interior-slope lower jets.

\smallskip\noindent
\underline{\textit{Step 5.}} Singular Hessians and unit slopes.

If a lower jet has positive-semidefinite singular spatial Hessian, the
pointwise scheme monotonicity $U_ m^{k+1}\geq U_ m^k$ passes to the
limit.  Testing this temporal monotonicity from the past gives
$\phi_ s\geq0=\mathcal G(H)$.  Finally, if $|a_ 0|=1$, convexity and the
$1$-Lipschitz bound imply
\begin{equation*}
\underline U(s_ 0,q_ 0+ta_ 0)=\underline U(s_ 0,q_ 0)+t
\qquad(t\geq0).
\end{equation*}
The lower-contact inequality along this ray gives
$a_ 0^THa_ 0\leq0$.  For $H\geq0$, this forces $Ha_ 0=0$, so $H$
is singular and the preceding argument applies.  Every subsequential
limit is therefore a strong convex viscosity solution of (44).  To
identify all subsequential limits, it remains to prove comparison in
this precise class.

\smallskip\noindent
\underline{\textit{Step 6.}} Whole-space localization for comparison.

For completeness, we prove comparison in the precise class produced by
(46)--(46a): the functions are continuous, convex and $1$-Lipschitz
in space, nondecreasing in time, have bounded cone corrections on
finite time intervals, and share a spatially uniform time modulus.
Let $u$ be a subsolution and $v$ a strong convex supersolution in
this class, with $u(0,\cdot)\leq v(0,\cdot)$, and suppose that
$u-v>0$ somewhere before a time $S$.

We need an interior maximum to which the doubled-variable lemma
applies.  The terminal penalty below creates a positive gap between
the time derivatives, the relative quadratic controls $x-y$, and the
last term makes the common spatial position bounded.
Choose the parameters in the order $\gamma$, then $\epsilon$, then
$\eta$.  Take $\gamma>0$ small enough that the violation survives
the terminal penalty, and then $\epsilon>0$ small enough for (48a)
to exclude the initial time.  Finally choose $\eta>0$ for spatial
coercivity; it will tend to zero with $\epsilon$ fixed.  Maximize
\begin{equation}
\begin{split}
\mathcal M_ {\epsilon,\eta}(t,x,y)
={}&u(t,x)-v(t,y)-\frac{|x-y|^2}{2\epsilon}
-\frac{\gamma}{S-t}\\
&-\eta\{\chi(x)+\chi(y)\},\qquad
\chi(x)=\sqrt{1+|x|^2},                                  
\end{split}
\tag{48}
\end{equation}
over $0\leq t<S$.  Since
\begin{equation*}
u(t,x)-v(t,y)\leq |x-y|+2C_ S,
\end{equation*}
the relative quadratic controls $x-y$, and the last term is
coercive in the common-translation direction.  Thus the positive
maximum is attained.  At $t=0$, initial ordering and the spatial
Lipschitz bound of $v$ give
\begin{equation}
u(0,x)-v(0,y)-\frac{|x-y|^2}{2\epsilon}
\leq |x-y|-\frac{|x-y|^2}{2\epsilon}\leq\frac\epsilon2.
\tag{48a}
\end{equation}
Taking $\epsilon<2\gamma/S$, the common time modulus keeps positive
maximizers a fixed distance from $t=0$.  The term
$-\gamma/(S-t)$ keeps them a fixed distance from $S$.

Before comparing Hessians, we remove the coercive term.  Although
the common center may then escape, translation by this center leaves
the relative quadratic unchanged.  Fix $\epsilon$, let
$\eta_n\downarrow0$, and let
$(t_ n,x_ n,y_ n)$ be maximizers.  The preceding cone estimate and
positivity of the maximum show that $d_ n=x_ n-y_ n$ is bounded.  Pass
to a subsequence with
\begin{equation*}
t_ n\to t_ *\in(0,S),\qquad d_ n\to d.
\end{equation*}
The common centres $c_ n=(x_ n+y_ n)/2$ may escape.  Translate both
functions by this same centre, and normalize only their vertical
constants:
\begin{equation}
\widehat u_ n(t,z)=u(t,z+c_ n)-u(t_ n,x_ n),\qquad
\widehat v_ n(t,z)=v(t,z+c_ n)-v(t_ n,y_ n).
\tag{48b}
\end{equation}
The spatial Lipschitz bound, common time modulus, and boundedness of
$d_ n$ make these families locally uniformly bounded and
equicontinuous.  After a subsequence they converge locally uniformly
to $\widehat u,\widehat v$.

Spatial translations and addition of constants preserve the equation.
The usual nearby-contact argument shows that $\widehat u$ is a
subsolution and $\widehat v$ is a strong convex supersolution.  Here
are the two boundary cases of that stability argument.  At a singular
upper test of $\widehat u$, add
$\sigma|z-z_ 0|^2$, pass to nearby contacts, and then let
$\sigma\downarrow0$.  At a positive-definite lower test of
$\widehat v$, nearby Hessians remain positive definite.  At a
singular positive-semidefinite lower test, temporal monotonicity,
which also passes to the limit, directly gives
$b\geq0=\mathcal G(Y)$.  Thus the strong formulation is closed under this
local limit.

Put $\xi_ n=d_ n/2$ and $\zeta_ n=-d_ n/2$.  Maximality in (48),
after the normalizations (48b), says, for every fixed $t,z,w$,
\begin{equation}
\begin{split}
&\widehat u_ n(t,z)-\widehat v_ n(t,w)
-\frac{|z-w|^2}{2\epsilon}-\frac{\gamma}{S-t}\\
&\quad\leq
-\frac{|\xi_ n-\zeta_ n|^2}{2\epsilon}
-\frac{\gamma}{S-t_ n}
+\eta_ n\{\chi(c_ n+z)+\chi(c_ n+w)
-\chi(x_ n)-\chi(y_ n)\}.                         
\end{split}
\tag{48c}
\end{equation}
Because $\chi$ is $1$-Lipschitz, the absolute value of the last
line, on compact $z,w$-sets, is at most
\begin{equation}
\eta_ n\{|z-\xi_ n|+|w-\zeta_ n|\}\longrightarrow0.
\tag{48d}
\end{equation}
Passing to the limit, while retaining the relative quadratic, shows
that $(t_ *,d/2,-d/2)$ is an attained global maximum of
\begin{equation}
\widehat u(t,z)-\widehat v(t,w)
-\frac{|z-w|^2}{2\epsilon}-\frac{\gamma}{S-t}.
\tag{48e}
\end{equation}
\smallskip\noindent
\underline{\textit{Step 7.}} Matrix order and comparison.

At the maximum in (48e), the relative quadratic is now the only
spatial penalty.  Its Hessian vanishes when both points move by the
same vector.  We use these directions in the doubled-variable matrix
inequality to obtain $X\leq Y$ with no localization error.

We use the doubled-variable lemma~\cite{CIL1992} in the following form.
Apply quadratic sup-convolution to $\widehat u$ and quadratic
inf-convolution to $\widehat v$.  The ordinary second-derivative test
at the regularized maximum, followed by passage to closed semijets, gives
\begin{equation}
(a,p,X)\in\overline{\mathcal P}^{2,+}\widehat u,\qquad
(b,p,Y)\in\overline{\mathcal P}^{2,-}\widehat v,\qquad
a-b=\frac{\gamma}{(S-t_ *)^2}>0,
\tag{48f}
\end{equation}
and, for every $\lambda>0$,
\begin{equation}
\begin{pmatrix}X&0\\0&-Y\end{pmatrix}
\leq A_ \epsilon+\lambda A_ \epsilon^2,\qquad
A_ \epsilon=\frac1\epsilon
\begin{pmatrix}I&-I\\-I&I\end{pmatrix}.
\tag{48g}
\end{equation}
The lower block bound in the same regularized second-derivative test
keeps the approximate Hessians bounded, so (48g) survives in the
closed semijets.  Both matrices on its right annihilate every common
direction $(z,z)$.  Testing (48g) on such a direction gives the
exact, error-free inequality
\begin{equation}
X\leq Y.
\tag{48h}
\end{equation}

Convexity of $\widehat u(t_ *,\cdot)$ gives $X\geq0$: apply each
approximating upper quadratic contact in opposite spatial directions
and use nonnegativity of the convex second difference.  Hence $Y\geq0$.
To apply the supersolution inequality to the closed lower jet, first
suppose $Y>0$.  The approximating Hessians are then eventually positive
definite, so the inequality passes to the limit by continuity.
If $Y$ is singular, temporal monotonicity gives $b\geq0=\mathcal G(Y)$:
at an ordinary lower contact, fixing the spatial point and approaching
from earlier times gives a nonnegative time derivative, and this
property passes to closed jets.  Since determinant is nondecreasing
in the Loewner order on $\mathbb S_ +^2$,
\begin{equation*}
a\leq \mathcal G(X)\leq \mathcal G(Y)\leq b,
\end{equation*}
contradicting the strictly positive time-derivative gap in (48f).
Comparison follows.  Applied to the upper and lower half-relaxed
limits, which have the same initial trace, it shows that they agree;
the compactness estimate then gives local uniform convergence of the
whole sequence.

\smallskip\noindent
\underline{\textit{Step 8.}} Scaling and the value at the origin.

To identify the normalization in (45), observe that the equation and
its initial datum are invariant under
\begin{equation*}
U(s,q)\mapsto a^{-1}U(a^2s,aq).
\end{equation*}
Uniqueness therefore gives
$U(s,q)=\sqrt s\,U(1,q/\sqrt s)$.  To obtain (45) for the prescribed
indices $k_m$, apply the convergence just proved with
$r_m=t_{k_m}^{(m)}$.  The two assumptions in (45) are precisely the
conditions required for this choice of $r_m$.  Formula (43) at
$(s,q)=(1,0)$ then gives the claimed limit $U(1,0)=\kappa$.

\end{proof}

\section{Completion of the Gaussian and empirical limits}\label{sec:conclusion}

We now choose the scales for which the two-scale comparison can be iterated
and the viscosity limit applied in the same sequence.  The proof keeps
track of three requirements: few omitted levels, a negligible sum of
comparison errors, and a vanishing time step after rescaling.

\begin{lemma}[The dyadic Gaussian limit]\label{lem:dyadic-gaussian-limit}
The dyadic Gaussian limit exists, and in the notation of Lemma~\ref{lem:viscosity-limit} it is
\begin{equation}
\lim_ {J\to\infty}
\frac{\Gamma_ {2^J}}{2^J\sqrt{J\log2}}=\kappa.
\tag{49}
\end{equation}

\end{lemma}

\begin{proof}
It suffices to approximate the Gaussian cost per unit area by the cell
iteration with error $o(\sqrt J)$.  Suppose that we retain $m$ groups
of $h$ levels, omit $D$ levels, and incur an error at most $e_J$ per
unit area at each replacement.  Put $t_J=mh\log2$.  The three
requirements are
\begin{equation*}
\frac{D}{J}\longrightarrow0,\qquad
e_J=o\!\left(\frac{h}{\sqrt J}\right)
\ \Longrightarrow\
m e_J=o(\sqrt J),\qquad
\frac{h}{t_J}\longrightarrow0.
\end{equation*}
The first allows us to restore omitted levels using (12), the second
controls the sum of the comparison errors, and the third permits the
viscosity limit.  We choose all parameters in the same sequence and
then carry out these reductions.

\smallskip\noindent
\underline{\textit{Step 1.}} Choose and align the retained bands.

Fix a large $J$, set
\begin{equation*}
h=\lfloor J^{3/4}\rfloor,\qquad b=\lceil(\log J)^2\rceil,\qquad
S=2^{b+1},
\end{equation*}
and retain consecutive groups of $h$ Haar levels separated by gaps of
$b$ omitted levels.  Omit also fewer than $h+b$ levels at the fine
end and $h+b$ levels at the coarse end.  The total number $D$ of
omitted levels satisfies
\begin{equation}
D\leq C\left(h+\frac{Jb}{h}\right)=o(J).
\tag{50}
\end{equation}
The coarse-end omission ensures that, after normalizing any retained
band $\eta_ {r,h}$ to unit outer blocks, the ambient square has side
\begin{equation}
R_ r=2^{J-r-h}\geq2^{h+b}.
\tag{50a}
\end{equation}
When it exists, the next retained detail starts at level $r+h+b$, whose
parent side is $2^{r+h+b+1}$.  Hence its smallest parent side in the
active-band normalization is exactly
\begin{equation*}
\frac{2^{r+h+b+1}}{2^{r+h}}=2^{b+1}=S.
\end{equation*}
Since $h\geq1$, (50a) also shows that $S$ divides $R_r$.  Thus the
ambient square is divisible by every auxiliary dyadic length used
below.  Orthogonality of the Haar projections makes the active band
independent of all bands still present.

\smallskip\noindent
\underline{\textit{Step 2.}} Make the comparison errors summable.

In Lemma~\ref{lem:two-scale-scalarization} choose
\begin{equation}
L=S^{1/8},\qquad B=S^{3/8},\qquad A=S^{5/8},
\qquad\varepsilon=S^{-1/32},
\tag{51}
\end{equation}
rounding to compatible dyadic integers.  The parameters are chosen in
the order
\[
J;\qquad h,b,S;\qquad L,B,A,\varepsilon.
\]
The omitted-level estimate is already (50).  To bound $e_J$ uniformly
throughout the iteration, we check both the geometry and the sizes of
the profiles in \eqref{eq:two-scale-scalarization}.  The geometric
slack required in (23a) is
\begin{equation*}
L+\frac{B^2}{A}\asymp S^{1/8}
\ll S^{11/32}=\varepsilon B.
\end{equation*}

The geometric loss satisfies $\delta\leq CS^{-1/32}$.
The concentration error is $O(\sqrt{\log S}\,S^{-1/16})$.
Since $S=2^{b+1}$ grows faster than any power of $J$, it remains to
bound the other factors polynomially in $J$, uniformly in the step.
If
$\Phi_ 0=|\cdot|$ and
$\Phi_ {k+1}=\mathcal H_ {h,L}\Phi_ k$, the construction behind (46)
gives
\begin{equation}
\Phi_ k(p)\leq\mathbb E|p+Z_ 1+\cdots+Z_ k|
\leq |p|+C\sqrt{kh};
\tag{52}
\end{equation}
hence the height parameter $M$ in (13) is $O(\sqrt J)$.
For the normalized retained band at level $r$, the canonical coarse flux
gives, uniformly in the step,
\begin{equation}
|Q_{R_r}|^{-1}\mathbb EJ_ {\mathcal H_ {h,L}\Phi_ k}^{Q_{R_r}}(\zeta)
\leq\mathbb E\fint_{Q_{R_r}}\bigl(|P|+C\sqrt J\bigr)\leq C\sqrt J.
\tag{52a}
\end{equation}
Inserting these bounds into \eqref{eq:two-scale-scalarization} gives
the required uniform estimate
\begin{equation}
e_ J\leq C\sqrt J\,S^{-1/32}
+C\sqrt{\log S}\,S^{-1/16}=o(h/\sqrt J).
\tag{52b}
\end{equation}
Since $m\leq J/h$, the second error budget is now verified:
\begin{equation*}
m e_J\leq\frac{J}{h}\,
o\!\left(\frac{h}{\sqrt J}\right)=o(\sqrt J).
\end{equation*}
Together with $\varepsilon L=S^{3/32}\to\infty$, the displayed slack
verifies the geometric conditions (23a).
All iterated profiles are $D_ 4$-invariant by Proposition~\ref{prop:cell-operator-properties}, so the reflected bounded-square proof applies.

\smallskip\noindent
\underline{\textit{Step 3.}} Remove the retained bands and evaluate the terminal cost.

Apply \eqref{eq:two-scale-scalarization}, successively, from the finest retained band to the coarsest.
There are $m\leq J/h$ applications, so their total error per unit
area is at most $me_ J=o(\sqrt J)$.  In dimension two the rescaling
from a band of outer side $2^{r+h}$ changes both total energy and area
by the same factor $2^{2(r+h)}$, so \eqref{eq:two-scale-scalarization} is an error per original unit
area with no additional scale factor.  Formula (50a) provides exact
boundary alignment at every application.  For the last band the same argument is used
with $\zeta=0$.  After the last replacement no random source remains.
For a divergence-free zero-normal field $v$ on a square,
\begin{equation*}
\int v_ i=\int\operatorname {div}(x_ i v)=0;
\end{equation*}
therefore Jensen's inequality, with equality for $v=0$, says that the
terminal energy is exactly
$|Q_{2^J}|\Phi_m(0)=2^{2J}\Phi_m(0)$.  Consequently, if
$\rho_{2^J}^{\rm ret}$ is the retained source in $Q_ {2^J}$,
\begin{equation}
2^{-2J}\mathbb E\|\rho_{2^J}^{\rm ret}\|_ {\rm KR}
=\Phi_ m(0)+o(\sqrt J).
\tag{53}
\end{equation}

\smallskip\noindent
\underline{\textit{Step 4.}} Apply the viscosity limit and restore the omitted bands.

The retained logarithmic time is
\begin{equation}
t_ J=mh\log2=(J-D)\log2=J\log2+o(J).
\tag{54}
\end{equation}
In particular,
\begin{equation*}
t_J\sim J\log2,\qquad
\frac{h}{t_J}\longrightarrow0,\qquad L=S^{1/8}\longrightarrow\infty.
\end{equation*}
This verifies the final error budget and precisely the joint limit regime
required by Lemma~\ref{lem:viscosity-limit}.

Apply Lemma~\ref{lem:viscosity-limit} with triangular row index $J$,
$h_J=h$, $L_J=L$, $r_J=t_J$, $F_k^{(J)}=\Phi_k$, and iteration index
$k_J=m$.  Then $t_{k_J}^{(J)}=t_J$, and (54) verifies the hypotheses
of (45).  Hence $\Phi_m(0)/\sqrt{t_J}\to\kappa$.  On the other hand,
(12) and (50) give
\begin{equation}
\mathbb E\|\rho_{2^J}-\rho_{2^J}^{\rm ret}\|_ {\rm KR}
\leq C2^{2J}\sqrt D=o(2^{2J}\sqrt J).
\tag{55}
\end{equation}
Combining (53)--(55),
\begin{equation}
\frac{\mathbb E\|\rho_{2^J}\|_ {\rm KR,Q_ {2^J}}}
{2^{2J}\sqrt{J\log2}}\longrightarrow\kappa.
\tag{56}
\end{equation}
Finally (9) is $o(2^{2J}\sqrt J)$, and
$\widehat\Gamma_ {2^J}=2^J\Gamma_ {2^J}$.  Equation (56) is therefore
equivalent to (49), which proves the lemma.

\end{proof}

\begin{samepage}
\begin{proof}[Proof of Theorem~\ref{thm:intro-main}]
The dyadic limit in (6) equals $\kappa$ by
Lemma~\ref{lem:dyadic-gaussian-limit}.  The approximation in
Lemma~\ref{lem:gaussianization-dyadic} therefore gives, for the full
sequence of integers $N$,
\begin{equation*}
\lim_ {N\to\infty}\frac1{\sqrt{N\log N}}
\mathbb E\sup_ {\Lip(f)\leq1}
\left|\sum_ {i=1}^N\left(f(X_ i)-\int f\,d\lambda\right)\right|
=\frac\kappa{\sqrt2}.
\end{equation*}
By Lemma~\ref{lem:viscosity-limit}, $\kappa=U(1,0)$.  The AKT bounds
recalled in the introduction show that this limit is finite and positive,
which proves every assertion of the theorem.

\end{proof}

\end{samepage}

\section*{Acknowledgement}

H.M. is supported in part by a Lady Davis Fellowship at the Technion.
The initial proof idea was developed by OpenAI's GPT-5.6 sol and
Rethlas proof agent.

\begingroup
\raggedright
\bibliographystyle{alpha}
\bibliography{biblio}
\endgroup

\end{document}
